\documentclass[11pt]{article}

\usepackage[margin=1in]{geometry}
\usepackage{amsmath,amssymb}
\usepackage{booktabs}
\usepackage[T1]{fontenc}
\usepackage[utf8]{inputenc}
\usepackage{lmodern}
\usepackage{microtype}
\usepackage{hyperref}
\usepackage{xurl}
\usepackage{enumitem}
\usepackage{longtable}
\usepackage{array}
\usepackage{xcolor}
\usepackage{pdflscape}

\hypersetup{
  colorlinks=true,
  linkcolor=blue,
  urlcolor=blue,
  citecolor=blue,
  pdftitle={The Current scripts-paper Pipeline of hetid},
  pdfsubject={Static source trace of the paper pipeline entrypoint}
}

\DeclareRobustCommand{\code}[1]{%
  \ifmmode\text{\ttfamily\detokenize{#1}}\else\nolinkurl{#1}\fi}
\DeclareRobustCommand{\file}[1]{\path{#1}}
\DeclareRobustCommand{\src}[2]{\file{#1}\,::\,\code{#2}}
\newcommand{\trans}{^{\top}}
\newcommand{\R}{\mathbb{R}}
\newcommand{\Covhat}{\widehat{\operatorname{Cov}}}
\newcommand{\Varhat}{\widehat{\operatorname{Var}}}
\newcommand{\E}{\operatorname{E}}
\newcommand{\Var}{\operatorname{Var}}
\newcommand{\Cov}{\operatorname{Cov}}
\newcommand{\TauP}{\textnormal{Target P}}
\newcommand{\TauS}{\textnormal{Target S}}

\setlist[itemize]{leftmargin=1.5em}
\setlist[enumerate]{leftmargin=1.7em}
\setlength{\LTpre}{0.5em}
\setlength{\LTpost}{0.5em}
\renewcommand{\arraystretch}{1.12}
\setlength{\emergencystretch}{1.5em}

\title{The Current \texttt{scripts-paper} Pipeline of \texttt{hetid}}
\author{}
\date{Source inspection: 1 August 2026, 08:59 EDT\\
Commit \texttt{c42cb4d273771871259986ef979eece3eee3abc2}}

\begin{document}

\maketitle

\begin{abstract}
The current paper pipeline is one ordered graph evaluated in a shared
environment. It has one unified bootstrap stage and one manifest-driven
artifact lifecycle. This static trace maps runner order, in-memory state,
scientific estimands, numerical certificates, inference rules, cache identity,
    the complete input and control interface, the exact default route,
    source-fixed global contracts, alternative code paths, and publication
    artifacts to the current R source. Runtime completion and
artifact freshness remain outside the trace. The account follows
\file{scripts-paper/run_pipeline.R} in source order and then records its
mathematical and typed-artifact contracts. The package layer under \file{R/}
owns the Lewbel moment and quadratic arithmetic. The paper layer under
\file{scripts-paper/} owns the application, numerical searches, inference,
diagnostics, figures, tables, reports, and artifact lifecycle. This review
inspected source without running the pipeline, initializing its caches, or
using generated output as evidence.
\end{abstract}

{\small\tableofcontents}

\section{What this static trace establishes}
\label{sec:contracts-summary}

The current source has the following main contracts:
\begin{itemize}
  \item The runner reads six environment inputs and no command-line flags. It
    also consumes four tracked plan or decision inputs: the mean-specification
    plan and the LAD, joint-GMM, and EGARCH-X decisions. The current LAD decision
    selects a dependency-gated stage. The current joint-GMM decision selects only
    replication checks, and every nondefault decision stops before substantive
    joint-GMM work. The EGARCH-X decision must match a fresh runtime gate, but no
    EGARCH-X producer exists in the runner. The reset entrypoint's
    \code{--keep-tracked} flag belongs to a separate maintenance command. Other
    editable global values are source-fixed code contracts, not launch-time
    options.
  \item Published mean specification B residualizes each
    stochastic-discount-factor (SDF) news principal-component (PC) score on
    the mean-control block $X$ and estimates the news-on-controls coefficient
    matrix $\beta_2^R$. Diagnostic mean specification A imposes
    $\beta_2^R=0$. The published expected-SDF and combined variance shares are
    set-valued.
  \item One unified bootstrap stage computes the published mean and
    log-variance results. Its primary resample-index family uses 10,000 circular
    moving-block-bootstrap (MBB) replications. The shared bootstrap frame has
    $T$ rows, and the primary family uses block length
    $\ell=\lceil1.5T^{1/3}\rceil$, seed 20260708, and
    Mersenne-Twister/Inversion/Rejection random-number-generator kinds. The
    volatility-sensitivity
    resample-index family uses block length $2\ell$.
    Diagnostic mean specification A reuses the primary resample-index family in
    the paired B-versus-A mean-specification comparison.
  \item At identification relaxation value $\tau=0$, the mean system has one
    authoritative mean news-coefficient point. Each estimator-specific analysis
    stage evaluates its log-variance estimator directly at that point. The
    tau-zero reporting gate either supplies a coefficient's bootstrap
    point/median-absolute-deviation (MAD) statistic and stars or records why
    they are unavailable.
  \item Positive relaxation values have side-specific endpoints. The published
    Target P confidence interval, whose target is pointwise coverage uniformly
    over the relevant scalar identified interval, appears only for a supported
    one- or two-sided status pattern that passes the endpoint regularity gate.
    Stored endpoint status has four values, including collector-injected
    \code{failed}; a finite value without \code{bounded} is not publishable.
  \item Poisson pseudo-maximum likelihood (PPML) and the Harvey estimator are
    primary log-variance estimators evaluated by the unified bootstrap stage.
    Least absolute deviations (LAD) is a dependency-gated extension outside that
    stage. The current source approves LAD with exactly \code{quantreg} 6.1; a
    missing or different version stops the run rather than skipping LAD.
    The exponential generalized autoregressive conditional
    heteroskedasticity model with exogenous regressors (EGARCH-X) has a decision
    router but no artifact producer in the current runner.
  \item PPML and the Harvey estimator receive the same primary selected-grid
    cap, scan-derived start pool, and fit-evaluation budget. Both second search
    audits use five scan-derived starts per coefficient and endpoint side, a
    40,000-fit-evaluation budget per displayed relaxation value, one fresh
    set-engine cache, and the same endpoint reconciliation. Only PPML's second
    search audit enlarges the selected-grid cap and changes its selection rule.
    Section~\ref{sec:ppml-harvey-search}
    states the full matched comparison.
  \item Bootstrap cache schema 4 and semantic cache freshness govern reuse of
    the unified bootstrap cache artifact. Separate mean and log-variance cache
    files are legacy names.
  \item The bounds-by-relaxation-value grid is tail-refined and nonuniform. The
    joint-null witness and joint-generalized-method-of-moments (GMM) replication
    products are diagnostics; they add no test or identifying restriction.
  \item The artifact manifest has 84 artifact records. Final
    conditional-route reconciliation covers the seven conditional records and
    does not certify the 77 required files.
\end{itemize}

\section{How to read this document}
\label{sec:reading}

The trace defines common abbreviations before the detailed terminology reference.
They include Adrian--Crump--Moench (ACM), Federal Reserve Economic Data (FRED),
ordinary least squares (OLS), and principal component analysis (PCA). Estimation
and inference terms include quasi-maximum-likelihood estimator (QMLE),
moving-block bootstrap (MBB), random-number generator (RNG),
heteroskedasticity-and-autocorrelation-consistent (HAC), and standard error (SE).
Artifact formats include comma-separated values (CSV), R data serialization
(RDS), Portable Document Format (PDF), and Scalable Vector Graphics (SVG).

\subsection{The two code layers}

The repository contains two distinct layers. The R package under \file{R/}
provides the reusable identification application programming interface (API).
It computes the centered moments,
builds the quadratic system, checks the constraints, and supplies bond-pricing
variance-bound kernels. The paper pipeline under \file{scripts-paper/} is a
self-contained analysis graph evaluated in a shared environment. It calls the
installed package for package-owned arithmetic and keeps paper-owned code under
\file{scripts-paper/support/}.

This distinction matters for code matching. A statement about the seven
identification moments should point to the package. A statement about
published mean specification B, the unified bootstrap, Target P inference, or a table belongs to
the paper layer. The authoritative project summary is \file{CLAUDE.md}; the
module catalog and production entrypoint appear in
\file{scripts-paper/README.md}.

\subsection{Scope of the trace}

The trace covers every direct top-level effect in
\file{scripts-paper/run_pipeline.R}, including its 47
\code{paper_source_once()} call sites. After prior statements succeed, the
runner reaches 46 call sites unconditionally. It conditionally sources only
\file{log_variance/estimators/lad/run_sets.R}. The runner sources
\file{log_variance/figures/fitted_volatility/run_lad.R} unconditionally; that
module guards its own body on the existence of the LAD result.

The document also covers production-reachable helper modules, package arithmetic,
the artifact manifest, and the static contract tests. Establishing current
pipeline completion requires runtime evidence because the runner preserves
required outputs until their producers overwrite them.

The word \emph{input} includes the launch environment, tracked configuration and
decision files, installed packages, external data, platform-derived values, and
the unified bootstrap cache artifact. The trace records a helper argument when a
production caller fixes it or when it selects a real branch. It does not call
every formal argument of every internal helper a user option. Tests and offline
helpers can pass such arguments directly, but the production runner cannot.

\subsection{Status words used in this document}
\label{sec:status-words}

The following status words make execution claims explicit:
\begin{itemize}
  \item \textbf{Always attempted} means the runner reaches the statement whenever
    every earlier statement has succeeded. It does not mean the statement cannot
    fail.
  \item \textbf{Selected path} means the current checked-in source and the default
    launch environment choose the guarded branch. A required package, file, or
    runtime check may still make the run stop.
  \item \textbf{Alternative path} means an implemented branch is not chosen by a
    plain invocation but can be chosen by the stated environment input or tracked
    source decision.
  \item \textbf{Source-selected} means a checked-in contract or decision record
    chooses between the selected path and an alternative path.
  \item \textbf{Launch-selected} means an environment input chooses between the
    selected path and an alternative path. The document states the path taken when
    the input is unset.
  \item \textbf{Runtime-selected outcome} means data, cache validation, or
    numerical checks choose the outcome during the run. Source inspection can
    state the complete decision rule but cannot assert the realized outcome.
  \item \textbf{Available but not called} means code implements the functionality,
    but \file{run_pipeline.R} has no path that invokes it under the current graph.
  \item \textbf{Not implemented} means a schema, decision value, or artifact
    placeholder reserves the functionality, but the substantive solver or
    producer does not exist.
  \item \textbf{Separate entrypoint} means another maintenance, validation, test,
    or package command or API call owns the functionality. Running the pipeline
    does not invoke that entrypoint.
\end{itemize}

Thus \emph{source-selected} never means ``observed to have completed.'' It
reports what the checked-in source asks the next invocation to do. A plain
invocation means \code{Rscript scripts-paper/run_pipeline.R} with all six
pipeline environment inputs unset.

\subsection{Four ledgers}

Four recurring questions organize the explanation:
\begin{enumerate}
  \item \textbf{Order:} What does the runner evaluate next?
  \item \textbf{State:} Which in-memory objects does that stage consume and
    produce?
  \item \textbf{Science:} Which estimand, numerical certificate, or inference
    target does the stage implement?
  \item \textbf{Publication:} Which artifact IDs can the stage write, and under
    what condition?
\end{enumerate}
Section~\ref{sec:runner-ledger} answers the first question literally.
Sections~\ref{sec:data}--\ref{sec:publication} answer the middle questions in
the same order. Sections~\ref{sec:artifacts} and~\ref{sec:validation} give the
complete artifact and validation crosswalks.

\subsection{Canonical terminology and code-name mapping}
\label{sec:terminology-rule}

This document applies one rule throughout: one semantic concept has one prose
name. Literal identifiers remain unchanged inside \code{code}, \file{paths},
and \emph{source references}. When the code uses several names for one concept,
the prose chooses one canonical term and records the source names in the
terminology reference in Section~\ref{sec:terminology}. Executable assignments,
function calls, schemas, and mathematical identities decide whether two names
denote the same concept. Code comments supply supporting evidence about intent.

Similar names do not imply identical concepts. A relaxation value changes the
scientific identified set. A numerical tolerance controls an algorithm at a
fixed relaxation value. The terminology reference records this distinction
and the other pairs that readers must keep separate.


\section{Runner overview}
\label{sec:runner}

\subsection{Execution model}

The runner sources modules and evaluates top-level effects in order. The
path layer, \src{scripts-paper/config/paths.R}{paper_source_once}, normalizes an
absolute source path. Its source-once protocol distinguishes an absent registry
binding, \code{FALSE} while loading, and \code{TRUE} after loading.
A second request for a loaded path returns without evaluation. A request for a
path already loading raises a circular-dependency error. A failed load removes
its registry binding. Successful loads use \code{base::sys.source()} in
\code{.GlobalEnv} by default.

This architecture makes the source order a dataflow contract. Modules leave
named objects in the shared environment. Later modules consume objects such as
\code{set_id_mean_eq}, \code{mean_eq_bounds_tau}, \code{log_var_eq_ppml},
\code{set_id_boot}, and \code{log_var_eq_set_boot}. The topology checks reject
production calls of the form \code{source(paper_path(...))}; the path bootstrap
at the start of the runner is the intentional exception.

\subsection{Dependency flow}

Pipeline stages consume objects produced by earlier stages. They pass these
objects between source modules within one shared R process.

\begin{figure}[htbp]
\centering
\small
\begin{tabular}{c}
\fbox{\parbox{0.86\textwidth}{\centering
configuration, artifact manifest, publishing-run guard, startup cleanup of conditional artifacts, ACM load}}\\[0.4em]
$\downarrow$\\[-0.2em]
\fbox{\parbox{0.86\textwidth}{\centering
SDF-news, expected-SDF, and lagged expected-SDF panels; quarterly real
consumption growth; yield volatility; nominal asset-return PCs;
three SDF-PC score blocks}}\\[0.4em]
$\downarrow$\\[-0.2em]
\fbox{\parbox{0.86\textwidth}{\centering
exogenous-news OLS benchmark, published mean specification B, diagnostic mean
specification A, mean identified set, variance shares}}\\[0.4em]
$\downarrow$\\[-0.2em]
\fbox{\parbox{0.86\textwidth}{\centering
two-step log-OLS benchmark, attained mean coefficient ranges, PPML, Harvey
estimator, diagnostics, and routes,
dependency-gated LAD}}\\[0.4em]
$\downarrow$\\[-0.2em]
\fbox{\parbox{0.86\textwidth}{\centering
unified bootstrap stage, paired B-versus-A mean-specification comparison,
inference tables}}\\[0.4em]
$\downarrow$\\[-0.2em]
\fbox{\parbox{0.86\textwidth}{\centering
figures, heteroskedasticity-diagnostics table, SDF approximation-error variance
bounds, quoted-number reconstruction and assertion,
descriptive report}}\\[0.4em]
$\downarrow$\\[-0.2em]
\fbox{\parbox{0.86\textwidth}{\centering
conditional-route state record and final LaTeX-sidecar cleanup}}
\end{tabular}
\caption{Main dependency flow in \file{scripts-paper/run_pipeline.R}. ACM means
Adrian--Crump--Moench; SDF means stochastic discount factor; PC means principal
component; OLS means ordinary least squares; PPML means Poisson
pseudo-maximum likelihood; and LAD means least absolute deviations. The runner
evaluates the stages from top to bottom in one shared R environment, so each
box consumes state produced above it.}
\label{fig:dependency-flow}
\end{figure}

Figure~\ref{fig:dependency-flow} shows why source order is a dataflow contract:
reordering stages may change object availability.

\subsection{Exact route of a plain invocation}
\label{sec:current-route}

Table~\ref{tab:current-route} is the short answer to ``what path does the code
take?'' It starts from a plain invocation, so all six environment inputs are
unset. It separates a path selected by source from a result that only the run can
determine. Every row is conditional on all earlier rows succeeding.

\begin{landscape}
\begingroup
\footnotesize
\begin{longtable}{@{}>{\raggedright\arraybackslash}p{0.20\linewidth}
  >{\raggedright\arraybackslash}p{0.34\linewidth}
  >{\raggedright\arraybackslash}p{0.38\linewidth}@{}}
\caption{Exact default route and alternatives in the current source. Every
middle-column entry is the \emph{selected path}. Every executable right-column
alternative is an \emph{alternative path}; the text instead says \emph{available
but not called}, \emph{not implemented}, or \emph{separate entrypoint} when one
of those exact boundaries applies. ``Plain
invocation'' means \code{Rscript scripts-paper/run_pipeline.R} from the package
root with the pipeline environment inputs unset. A hard stop prevents every
later row from running.
\label{tab:current-route}}\\
\toprule
Stage reachability and selection basis & Selected path under a plain invocation & Alternative path or unavailable/separate functionality\\
\midrule
\endfirsthead
\toprule
Stage reachability and selection basis & Selected path under a plain invocation & Alternative path or unavailable/separate functionality\\
\midrule
\endhead
Working directory: always attempted &
Require the current directory to contain \file{DESCRIPTION} with package name
\code{hetid} and the \file{scripts-paper/} directory. &
Any other current directory stops before configuration or output handling. The
runner accepts no directory argument.\\
Bootstrap size and publishing guard: launch-selected &
Use 10,000 replications. The publishing guard passes without an acknowledgement. &
\code{HETID_BOOT_REPS} selects another whole number of at least 2. Any value other
than 10,000 stops unless \code{HETID_ALLOW_DRAFT_RUN} is nonempty. The
acknowledgement does not choose the number.\\
Bootstrap workers: launch-selected &
Use detected logical cores minus two on macOS or minus one elsewhere, with a
one-worker minimum. &
\code{HETID_BOOT_CORES} selects a whole number of at least 1. It changes
scheduling, not the prebuilt resample indices or semantic cache identity.\\
Bootstrap mode: launch-selected; runtime-selected outcome &
Request \code{reuse}. A present cache is reused only if every validation check
passes; otherwise the stage reports \code{fallback-rerun}, recomputes the whole
payload, and replaces the cache transactionally. On a clean state with no cache
file, this rule selects \code{fallback-rerun}. &
\code{HETID_BOOT_MODE=rerun} skips the cache read, recomputes the whole payload,
and reports \code{rerun}. No mode can reuse only one result branch.\\
Quadratic equivalence assertion: launch-selected &
Do not run the extra legacy-versus-generalized quadratic-builder assertion. &
Any nonempty \code{HETID_ASSERT_EQUIV} runs it at every reached quadratic-system
build. A mismatch stops; a match does not change the scientific object.\\
Artifact startup: always attempted &
Create manifest parent directories, then remove and audit every conditional LAD
and EGARCH artifact before either route is evaluated. Preserve regular artifacts
until their producers overwrite them. &
The runner exposes no skip-cleanup or full-reset branch. Full reset is a separate
command described in Table~\ref{tab:outside-runner}.\\
Quarterly ACM input: source-selected; runtime-selected outcome &
Request quarterly yields and term premia for months 1--120 with
\code{source="auto"}, \code{auto_download=FALSE}, and incomplete quarters retained.
The package resolves an already available valid source or stops. &
A source edit can select the package's GitHub or New York Fed monthly source.
The direct package API also has download controls; the paper runner neither
exposes them nor downloads this quarterly input.\\
Daily ACM input: source-selected; runtime-selected outcome &
Request daily yields for months 1--120 with \code{source="auto"} and
\code{auto_download=TRUE}. Use a valid cached asset when available or invoke the
package download path. &
A source edit can change the data contract. The New York Fed source cannot
supply the daily frequency, so it is not an alternative for this request.\\
FRED transport: launch-selected &
With \code{FRED_API_KEY} unset, fetch \code{PCECC96} through the public CSV
endpoint using the checked-in curl retry and timeout policy. &
A nonempty \code{FRED_API_KEY} selects the official JSON API and its separate
retry policy. Both transports feed the same consumption-series construction.\\
Data and mean specifications: source-selected &
Build all data series. Compute mean specification B first and publish it; compute
mean specification A second for diagnostics. Use B as the only volatility-panel
specification. &
A tracked \code{PAPER_SPEC_PLAN} edit may change the A/B order or use one mean
entry, subject to its assertions. The first mean entry becomes published, and
the sole volatility entry must equal it.\\
Core estimators and diagnostics: always attempted &
Run the two-step log-OLS benchmark, PPML, the Harvey estimator, the joint-null
witness diagnostic, the joint-GMM diagnostic, and the residual-dynamics gate. &
These are not run/skip options. Numerical status and diagnostic verdicts are
runtime-selected outcomes and may stop a later validation or publication stage.\\
Joint-GMM branch: source-selected &
Use the checked-in no-answer-default record and run only graph, algebra, and
software replication checks. The diagnostic does not gate later stages. &
The decision schema represents nondefault moment blocks, another intercept target,
stage C, and a tolerance grid. The current driver deliberately stops with
\code{unwired_joint_branch} on every nondefault record. Those branches are
available but not called by this runner.\\
Residual dynamics and EGARCH-X: source-selected; runtime-selected outcome &
Run the residual-dynamics gate, then require its fresh record to match the
checked-in EGARCH-X decision record. That record binds a non-rejecting verdict
and both decisions \code{not_asked}; a match closes the dynamic workstream and
requires all four conditional EGARCH artifacts to remain absent. &
A fresh gate mismatch stops rather than changing the committed decision. The
router defines rejecting-gate approval and dependency branches, but the current
runner sets no dependency variables and calls no EGARCH-X estimator or artifact
producer. Even an approval record cannot produce estimates in this graph.\\
LAD: source-selected; runtime-selected outcome &
The checked-in decision is \code{approved} for exactly \code{quantreg} 6.1. If
that version is installed, source and run LAD. If it is absent or different,
stop. The current source never silently converts this approval into ``skip LAD.'' &
An existing record with \code{unanswered} or \code{declined} selects the
alternative path that skips LAD without creating a registry entry. Deleting the
record does not select that path: the runner constructs it through
\code{paper_path()}, whose required-path check stops first. These are tracked
decision alternatives, not launch-time environment choices.\\
Unified bootstrap: always attempted after the gates &
Construct both 10,000-draw resample-index families. Run or reuse published mean
specification B plus PPML and Harvey result draws. Exclude LAD. &
Replication count, worker count, and cache mode change as described above. The
estimator capability registry, not a runtime estimator flag, fixes PPML and
Harvey as the primary bootstrap estimators.\\
Mean-specification comparison: source-selected &
Use the primary resample-index family to compute paired diagnostic draws for
mean specification A and write rows for published B and diagnostic A. &
If the mean plan has one entry, compute no alternative draws but still write the
required comparison artifact with the published specification.\\
Publication and final state: always attempted, with conditional LAD content &
Attempt every regular publication stage. Include LAD pages and figures only if
LAD ran. The fitted-volatility sweep first filters configured relaxation values
against the runtime boundedness frontier. If none survives, the sweep stops. If
one survives, it writes only the two estimator-specific scalar paths. If more
than one survives, it also writes the eight aggregate paths. It therefore may
omit manifest-required paths whenever it does not stop.
Record that the EGARCH-X producer did not run, write the conditional-route state
record, and remove LaTeX sidecars. &
No current branch publishes EGARCH-X artifacts, and no current final gate checks
all regular artifacts. A failure in any writer,
\code{latexmk}, route reconciliation, or final cleanup stops the run; later
success cannot be inferred from earlier artifacts.\\
\bottomrule
\end{longtable}
\endgroup
\end{landscape}

The table deliberately does not label cache reuse, LAD completion, external-data
resolution, or a diagnostic verdict as a source fact. Those are runtime-selected
outcomes. It labels B, A, the no-answer joint-GMM branch, and the approved LAD
decision as source-selected because checked-in code fixes them.

\subsection{Exact top-level ledger}
\label{sec:runner-ledger}

The runner follows a literal top-level order. At top level, a definition module
only supplies functions or constants. Paths in the source-or-effect column are relative to
\file{scripts-paper/}.
Every ledger row has status \emph{always attempted} after earlier rows succeed,
except lines 112--114, whose status is \emph{selected path} and source-selected, subject to
the LAD package gate. A row that always sources a file can still contain a
guarded body; the six driver-level body guards are listed after the table. This rule assigns an explicit
status to every row without repeating the same words 46 times.

\small
\begin{longtable}{@{}>{\raggedleft\arraybackslash}p{0.07\textwidth}
  >{\raggedright\arraybackslash}p{0.37\textwidth}
  >{\raggedright\arraybackslash}p{0.49\textwidth}@{}}
\caption{Top-level runner ledger for \file{scripts-paper/run_pipeline.R} at the
commit printed on the title page. ``Lines'' gives source-line positions in that
file; paths in ``Source or effect'' are relative to \file{scripts-paper/}. ACM
means Adrian--Crump--Moench, and API means application programming interface. PC
means principal component. OLS means ordinary least squares. PPML means Poisson
pseudo-maximum likelihood. QMLE means quasi-maximum-likelihood estimator. HAC
means heteroskedasticity-and-autocorrelation consistent. GMM means generalized
method of moments. EGARCH-X means exponential generalized
autoregressive conditional heteroskedasticity with exogenous regressors, and LAD
means least absolute deviations. SDF means stochastic discount factor. CSV means
comma-separated values. RDS means R data serialization. TeX is the LaTeX source
format. PDF means Portable Document Format. SVG means Scalable Vector Graphics.
The ledger's main result is the exact evaluation
order and the principal state or artifacts produced at each stage. Unless the
lines 112--114 row or a guarded body says otherwise, every row is always
attempted after prior success.
\label{tab:runner-ledger}}\\
\toprule
Lines & Source or effect & Role and principal product\\
\midrule
\endfirsthead
\toprule
Lines & Source or effect & Role and principal product\\
\midrule
\endhead
5--8 & \file{config/paths.R} via base \code{source()} &
Validate the package root and define paths, runtime controls, and the source-once registry.\\
9 & \file{config/artifacts.R} &
Load the artifact manifest, artifact query API, lifecycle, typed writers, LaTeX publication registry,
and reset definitions.\\
10 & \file{config/analysis.R} &
Load analysis axes, the log-variance estimator registry, reporting policy, bootstrap controls, and
\code{PAPER_SPEC_PLAN}.\\
16--24 & Publishing-run guard &
Stop before directory creation if \code{boot_reps != 10000} and
\code{HETID_ALLOW_DRAFT_RUN} is empty.\\
25 & \file{support/data/acm_inputs.R} &
Define the validated quarterly ACM loader.\\
26--41 & Directory creation and startup cleanup of conditional artifacts &
Create every artifact manifest parent and delete all conditional LAD and EGARCH paths,
recording audits of the startup cleanup of conditional artifacts.\\
42 & \code{paper_load_quarterly_acm(all_mats)} &
Load validated quarterly yields and term premia for months 1--120.\\
45 & \file{data_preparation/fred_download_patch.R} &
Install the retrying FRED method when \code{quantmod} is available.\\
\midrule
48 & \file{data_preparation/build_sdf_series.R} &
Produce \code{sdf_news}, \code{expected_sdf}, and \code{lag_expected_sdf}.\\
49 & \file{data_preparation/build_consumption_growth.R} &
Produce quarterly real consumption growth \code{gr1_pcecc96}.\\
50 & \file{data_preparation/build_yield_volatility.R} &
Produce the daily-change-based quarterly yield-volatility panel, including the
five-year-yield volatility instrument \code{y60_vol}.\\
51 & \file{data_preparation/build_asset_return_pcs.R} &
Load and quarter-align bundled nominal asset-return PCs as
\code{l.pc1}--\code{l.pc4}.\\
52 & \file{data_preparation/build_sdf_pcs.R} &
Produce the three SDF-PC score blocks.\\
53 & \file{mean_equation/fit_ols.R} &
Fit the exogenous-news OLS benchmark \code{ols_mean_eq}.\\
54 & \file{mean_equation/estimate_identified_set.R} &
Estimate published mean specification B and diagnostic mean specification A;
bind published mean specification B to \code{set_id_mean_eq}.\\
55 & \file{mean_equation/variance_shares/compute_variance_shares.R} &
Compute \code{var_share} for expected-SDF, SDF-news, and combined blocks.\\
56 & \file{mean_equation/variance_shares/render_variance_share_table.R} &
Publish the variance-share fragment, standalone TeX, and PDF.\\
57 & \file{log_variance/estimators/log_ols/run.R} &
Freeze the shared log-variance sample, build \code{log_var_eq}, and initialize the
bounds-rendering registry.\\
58 & \file{mean_equation/inference/compute_bounds_by_tau.R} &
Build \code{mean_eq_bounds_tau} and write the attained mean coefficient-range figure.\\
61 & \file{log_variance/estimators/ppml/run_sets.R} &
Map the PPML quasi-Poisson log-link QMLE over the mean identified set and run its
grid-expanded second search audit with an 8,000-row selected-grid cap.\\
63 & \file{log_variance/inference/standard_error_estimators.R} &
Define shared analytic covariance helpers.\\
66 & \file{log_variance/estimators/ppml/standard_errors.R} &
Attach PPML covariance variants; the paper selects four-lag HAC.\\
69 & \file{log_variance/estimators/harvey/run.R} &
Map the Harvey Gaussian multiplicative-variance QMLE over the mean identified set
and run its fixed-grid second search audit with a 4,000-row selected-grid cap; attach
Harvey estimator covariance variants.\\
\midrule
73 & \file{log_variance/diagnostics/joint_null/distance_objective.R} &
Define the joint-null objective.\\
74 & \file{log_variance/diagnostics/joint_null/search.R} &
Define its scan, multistart, and polish routines.\\
75 & \file{log_variance/diagnostics/joint_null/stability.R} &
Define crossing-geometry and witness-stability checks.\\
76 & \file{log_variance/diagnostics/joint_null/run.R} &
Run the joint-null witness diagnostic and write its CSV/RDS pair.\\
79 & \file{log_variance/diagnostics/joint_gmm/run.R} &
Validate the decision record and run the currently wired joint-GMM replication diagnostic.\\
86 & \file{log_variance/diagnostics/dynamics/gate_core.R} &
Define the residual-dynamics gate and gate-record builders.\\
87 & \file{log_variance/diagnostics/dynamics/run_gate.R} &
Evaluate the residual-dynamics gate using its lag-4 Ljung--Box rule and write
the residual-dynamics gate record.\\
93 & \file{log_variance/extensions/egarch/decision_core.R} &
Define the pinned EGARCH-X decision and routing validation.\\
94 & \file{log_variance/extensions/egarch/run_route.R} &
Validate the residual-dynamics gate record, route every branch, and rewrite the
terminal EGARCH status record.\\
101 & \file{log_variance/estimators/lad/dependency_gate.R} &
Define the LAD decision and exact dependency-version gate.\\
102--111 & LAD availability and decision read &
Construct \code{lad_gate}; approved but absent or version-mismatched \code{quantreg} is fatal.\\
112--114 & \file{log_variance/estimators/lad/run_sets.R} &
Selected path, source-selected: conditionally build and register LAD and its closure
diagnostics after exact \code{quantreg} 6.1 passes the package gate.\\
\midrule
118 & \file{inference/run_bootstrap_stage.R} &
Run or reuse the unified bootstrap stage; expose \code{set_id_boot},
\code{log_var_eq_set_boot}, and \code{bootstrap_primary_family}.\\
121 & \file{mean_equation/inference/spec_comparison.R} &
Bootstrap diagnostic mean specification A on the exact primary resample-index
family used for published mean specification B.\\
129 & \file{log_variance/tables/render_combined_inference_table.R} &
Publish the combined mean-equation and PPML inference table triple.\\
131 & \file{log_variance/tables/render_estimator_pages.R} &
Publish PPML and two-step log-OLS benchmark pages, then Harvey estimator and LAD
pages when their results exist.\\
134 & \file{log_variance/figures/render_bounds_by_tau.R} &
Render a bounds-by-relaxation-value figure for each bounds-rendering registry entry.\\
136 & \file{log_variance/figures/fitted_volatility/run.R} &
Render baseline PPML and Harvey estimator fitted
conditional-residual-volatility envelopes.\\
139--141 & \file{log_variance/figures/fitted_volatility/run_tau_sweep.R} &
Filter configured relaxation values by the runtime boundedness frontier; render
PPML and Harvey scalar panels for surviving values and render combined linear,
log, and normalized panels only when more than one value survives. Stop before
rendering if no configured value survives.\\
144 & \file{log_variance/figures/fitted_volatility/run_lad.R} &
Always source the module; render the LAD fitted conditional-median
absolute-residual envelope only when \code{log_var_eq_lad} exists.\\
145 & \file{mean_equation/figures/prepare_region_geometry.R} &
Prepare rescaling, refined boxes, grids, and envelopes for mean-region figures.\\
146 & \file{mean_equation/figures/render_projections.R} &
Render projections with and without the OLS benchmark.\\
147 & \file{mean_equation/figures/render_region_3d.R} &
Require all four configured relaxation values to lie below the runtime
boundedness frontier; then render 16 region figures across relaxation values,
units, and OLS modes.\\
148--153 & \file{mean_equation/diagnostics/heteroskedasticity/render_table.R} &
Run the heteroskedasticity diagnostic battery and publish its three-panel table
triple.\\
156 & \file{variance_bounds/compute_bounds.R} &
Compute SDF-news and expected-SDF approximation-error variance bounds for 39
quarterly maturities.\\
159 & \file{variance_bounds/quoted/run.R} &
Regenerate and assert quoted SDF approximation-error variance-bound numbers
before the bound renderers.\\
160 & \file{variance_bounds/figures/render_bounds.R} &
Write the SDF approximation-error variance-bound SVG.\\
161 & \file{variance_bounds/tables/render_summary_table.R} &
Publish the variance-bound summary triple.\\
162 & \file{reports/build_descriptive_statistics.R} &
Publish descriptive tables, a two-page figure PDF, wrapper TeX, and report PDF.\\
\midrule
165--170 & \code{build_conditional_route_status()} &
Reconcile route decisions, conditional requests, producer execution, and
conditional-artifact presence; carry audits of the startup cleanup of
conditional artifacts for provenance;
the runner passes
\code{egarch_producer_ran = FALSE}.\\
171--175 & \code{paper_write_exact_rds()} &
Write the required \code{conditional_route_status} conditional-route state record.\\
176 & \code{clean_latex_sidecars(out_dir)} &
Remove configured LaTeX sidecars recursively; this final action may still fail.\\
\bottomrule
\end{longtable}
\normalsize

Six always or conditionally sourced drivers guard their top-level bodies. The
joint-null driver requires \code{log_var_eq}, \code{set_id_mean_eq}, and
\code{mean_eq_bounds_tau}. The joint-GMM and residual-dynamics drivers each
require \code{log_var_eq} and \code{set_id_mean_eq}. The conditionally sourced
LAD driver again requires all three upstream objects. Every one of these
\code{exists()} predicates is true in runner order if earlier stages succeeded.
The EGARCH-X driver requires the fresh residual-dynamics RDS to exist; the
immediately preceding gate either writes it or stops, so that predicate is also
true when reached. The always sourced LAD fitted-scale driver is different: its
body runs only if the dependency-gated LAD result exists. These guards allow the
same files to be sourced by tests or offline tools without running their pipeline
bodies; they do not create hidden production choices.

Table~\ref{tab:runner-ledger} ends with conditional-route reconciliation and
LaTeX-sidecar cleanup. Startup cleanup of conditional artifacts and full
pipeline-state reset are different operations with different scopes.

\section{Configuration, preflight, and lifecycle}
\label{sec:preflight}

\subsection{Analysis contract}

The main scientific authority is
\src{scripts-paper/config/analysis_contract.R}{PAPER_ANALYSIS_CONTRACT}.
This authority fixes the axes and policies most likely to affect a reader's
interpretation.

\begin{table}[htbp]
\centering
\small
\caption{Current source-fixed analysis and inference contract. SDF means
stochastic discount factor; PC means principal component; PCA means principal
component analysis; MBB means moving-block bootstrap; SE means standard error;
HAC means heteroskedasticity-and-autocorrelation consistent; and PPML means
Poisson pseudo-maximum likelihood. Target P is the positive-relaxation-value
confidence-interval procedure. Its inference error rate is
$\alpha_{\mathrm{inf}}$. The default bootstrap has $B$ replications and a
shared frame of $T$ rows; $\ell$ is the primary MBB block length. The table
states source-fixed configuration; pipeline outcomes require runtime evidence.
\label{tab:contract}}
\begin{tabular}{@{}>{\raggedright\arraybackslash}p{0.31\textwidth}
  >{\raggedright\arraybackslash}p{0.59\textwidth}@{}}
\toprule
Field & Current value or rule\\
\midrule
Three SDF-PC score blocks & Three scores
in each block; each PCA centers and scales its input.\\
Nominal asset-return PCs & Four; downstream samples center them without scaling.\\
Analysis window & 1962 Q1 through 2025 Q4.\\
Mean specification plan & \code{c("B", "A")}; published mean specification B
is first and diagnostic mean specification A is second.\\
Mean specification used by log-variance estimators & Published mean
specification B only.\\
Relaxation-value policy & Baseline 0.05; display/projection 0.05, 0.10, 0.20; cap 0.99; sweep
step 0.005; fitted-volatility sweep over relaxation values through 0.40.\\
Inference policy & Version 2.0.0; $\alpha_{\mathrm{inf}}=0.10$; minimum usable share 0.50;
stability share 0.85; \TauP{} location tolerance $10^{-4}$.\\
Bootstrap & Default $B=10{,}000$; seed 20260708; mode \code{reuse} or \code{rerun}.\\
Primary resample-index family & Circular MBB with block length
$\ell=\lceil1.5T^{1/3}\rceil$.\\
Volatility-sensitivity resample-index family & Circular MBB with block length
$2\ell$.\\
Reporting SEs & PPML and the Harvey estimator both select Bartlett HAC with four lags.\\
\bottomrule
\end{tabular}
\end{table}

Table~\ref{tab:contract} shows that the volatility-sensitivity resample-index
family doubles the primary block length. PPML and the Harvey estimator retain
the same reporting covariance.

The maturity axes come from package constants. The quarterly SDF grid uses
months 3--117. The runner loads quarterly ACM yields and term premia for the
full 1--120 month grid before any data-construction source module. The
five-year-yield volatility instrument is \code{y60_vol}, the demeaned realized
quarterly volatility of the five-year yield.

\subsection{Complete input and prerequisite inventory}
\label{sec:input-inventory}

Table~\ref{tab:input-inventory} lists what can affect a production invocation.
It distinguishes an input consumed by the runner from a tool or file that merely
exists elsewhere in the repository.

\begingroup
\footnotesize
\begin{longtable}{@{}>{\raggedright\arraybackslash}p{0.22\textwidth}
  >{\raggedright\arraybackslash}p{0.35\textwidth}
  >{\raggedright\arraybackslash}p{0.35\textwidth}@{}}
\caption{Complete production-input inventory. Direct R namespaces are named so
that a package used only by dormant or test functionality is not mistaken for a
current prerequisite.
\label{tab:input-inventory}}\\
\toprule
Input class & Exact current input and selection & Use, alternative, and failure behavior\\
\midrule
\endfirsthead
\toprule
Input class & Exact current input and selection & Use, alternative, and failure behavior\\
\midrule
\endhead
Command and current directory &
The sole production command is \code{Rscript scripts-paper/run_pipeline.R}; it
must start at the \code{hetid} package root. The runner parses no command-line
arguments. &
\file{config/paths.R} reads \file{DESCRIPTION}, fixes absolute source paths and
the repository-relative output root, and stops on any other directory or package
name.\\
Checked-in source and records &
The runner consumes the \file{scripts-paper/} source tree, its artifact manifest,
the analysis and reporting contracts, \code{PAPER_SPEC_PLAN}, the joint-GMM and
EGARCH-X decision records, and \file{config/decisions/lad.dcf}. &
These files select the source-fixed behavior documented in
Tables~\ref{tab:source-controls} and~\ref{tab:current-route}. Editing one changes
the code or tracked decision, not a launch-time option; relevant files and whole
control objects also enter cache identity.\\
R and installed \code{hetid} &
The package declares R 4.5.0 or later. The runner calls the installed
\code{hetid} namespace for constants, ACM extraction, date normalization, SDF
construction, and package-owned identification arithmetic. &
The source checkout supplies the paper scripts, but it does not replace the
installed package namespace. Missing or incompatible functions stop at their
first call. The installed \code{hetid} version enters bootstrap runtime
provenance.\\
Direct namespaces on the regular path &
The regular production stages directly call \code{dplyr}, \code{ggplot2},
\code{knitr}, \code{nloptr}, \code{purrr}, \code{sandwich}, \code{skedastic},
\code{svglite}, \code{tibble}, \code{tidyquant}, \code{tidyr}, and \code{tsibble},
plus base or recommended R namespaces such as \code{parallel}, \code{stats},
\code{tools}, and \code{utils}. \code{tidyquant} reaches \code{quantmod} for FRED. &
There is no production lockfile or common preflight that checks all versions.
A missing namespace stops when its selected stage calls it. \code{nloptr} and
\code{sandwich} versions, together with R, platform, BLAS, and LAPACK details,
enter bootstrap runtime provenance.\\
Conditional namespaces and executables &
The JSON FRED route directly calls \code{jsonlite}; the selected LAD route calls
\code{quantreg} 6.1. The default FRED CSV route requires the command-line
\code{curl}. Table and report producers require \code{latexmk} and its LaTeX
engine. &
\code{jsonlite} is needed only when a nonempty \code{FRED_API_KEY} selects the API.
The current LAD approval makes exact \code{quantreg} 6.1 a run requirement rather
than an alternative fallback. A LaTeX failure stops its producer. \code{rugarch} is
named in the EGARCH-X decision protocol but is never loaded by the current
runner.\\
Quarterly ACM data &
Quarterly yields and term premia for months 1--120, with incomplete quarters
retained, \code{source="auto"}, and \code{auto_download=FALSE}. &
The package may resolve a valid existing GitHub-cache or bundled monthly asset.
The source-fixed paper request does not download this input. Missing or invalid
data, dates, columns, or maturity coverage stop before data construction.\\
Daily ACM data &
Daily yields for months 1--120 with \code{source="auto"} and
\code{auto_download=TRUE}. &
The package uses a valid cache or its digest-checked GitHub download path. These
data produce quarterly realized yield volatility and \code{y60_vol}. The New
York Fed source cannot satisfy a daily request.\\
FRED data &
\code{PCECC96} from 1947-01-01 through 2026-06-19. The default transport is the
public CSV endpoint; a nonempty API key selects JSON. &
The series produces quarterly real consumption growth. Transport failure or an
unusable series stops the data stage. A data change propagates through every
estimator and changes the bootstrap input hash.\\
Bundled nominal asset-return PCs &
The installed package data object \code{hetid::variables}, source columns
\code{pc1}--\code{pc4}, normalized to quarter end and shifted forward one quarter. &
The runner offers no alternate nominal asset-return-PC file. A source-contract
edit can change column names or dimension and then changes the log-variance frame
and bootstrap input.\\
Launch environment &
Exactly six variables are read: \code{HETID_BOOT_REPS},
\code{HETID_ALLOW_DRAFT_RUN}, \code{HETID_BOOT_CORES}, \code{HETID_BOOT_MODE},
\code{FRED_API_KEY}, and \code{HETID_ASSERT_EQUIV}. &
Their defaults, allowed values, and downstream effects appear in
Table~\ref{tab:user-controls}. No other production path rooted at the runner
reads a user environment variable.\\
Unified bootstrap cache artifact &
When mode resolves to \code{reuse}, the stage may read the one schema-4 cache at
the manifested bootstrap path. A clean state supplies no such input. &
Only a complete semantic and structural validation permits reuse. Missing,
unreadable, malformed, or stale content selects a complete
\code{fallback-rerun}; \code{rerun} mode does not read it. No generated result
outside this cache selects an estimator or scientific branch.\\
Preexisting output artifacts &
Conditional LAD and EGARCH artifacts are deleted before routing. Regular
artifacts are left in place until a producer replaces them. &
Regular output presence is not a branch input and does not prove that the
current invocation reached its producer. Conditional-route reconciliation
checks conditional presence only; it is not a completeness check for regular
artifacts.\\
Platform and derived provenance &
Logical-core detection and the operating system set the default worker count.
R version, platform, BLAS, LAPACK, selected package versions, data and design
hashes, code hashes, control objects, and resample-index and RNG hashes are
derived during bootstrap provenance construction. Git HEAD is read when
available for diagnostic provenance, and elapsed time is recorded. &
These are runtime inputs or metadata, not supported choices. A semantic
provenance mismatch prevents cache reuse. Missing Git metadata yields an
unavailable provenance stamp rather than selecting another scientific path;
clock time changes duration metadata, not estimates.\\
Writable filesystem and graphics/LaTeX runtime &
The output tree, temporary files, SVG/PDF devices, and LaTeX subprocesses must be
usable. &
I/O, device, compilation, readback, or final sidecar-cleanup failure stops the
runner. These facilities enable selected writers; they do not choose a model.\\
\bottomrule
\end{longtable}
\endgroup

\subsection{User controls and source-fixed contracts}
\label{sec:user-controls}

The runner and its production configuration consume ten top-level behavior
inputs: six environment inputs and four tracked plan or decision inputs. The
tracked inputs are the mean-specification plan and the LAD, joint-GMM, and
EGARCH-X decisions. The separate reset command exposes one maintenance flag.
Table~\ref{tab:user-controls} lists all eleven without implying that the reset
flag belongs to the runner. It separates the value
the user supplies from the consequence, so changing execution cost cannot be
mistaken for changing the scientific design.

\begingroup
\footnotesize
\begin{longtable}{@{}>{\raggedright\arraybackslash}p{0.24\textwidth}
  >{\raggedright\arraybackslash}p{0.23\textwidth}
  >{\raggedright\arraybackslash}p{0.45\textwidth}@{}}
\caption{Complete top-level behavior-input interface: ten pipeline inputs plus
the separate reset command's one maintenance flag. Default means the behavior when
the user supplies no override. A source edit is a tracked code change, not a
launch-time option.
\label{tab:user-controls}}\\
\toprule
Control and change route & Default and allowed values & Immediate and downstream effect\\
\midrule
\endfirsthead
\toprule
Control and change route & Default and allowed values & Immediate and downstream effect\\
\midrule
\endhead
Unified-bootstrap replication count:
\code{HETID_BOOT_REPS} & Unset or empty selects 10,000. An override must be a
whole number of at least 2. A value other than 10,000 also requires the
draft-run acknowledgement. & Sets the number of draws in both the primary and
volatility-sensitivity resample-index families. It therefore changes mean and
PPML/Harvey bootstrap draws, endpoint gates, Target P intervals, tau-zero
p-values and stars, volatility sensitivity, and the draw family available to
the paired mean-specification comparison. The count and family hashes enter
semantic cache provenance, so changing the count invalidates unified-bootstrap
reuse.\\
Bootstrap worker count: \code{HETID_BOOT_CORES} & An override must be a whole
number of at least 1. By default, the runner uses detected logical cores minus
two on macOS and minus one on other systems, with a one-worker minimum; failed
detection also gives one worker. & Controls parallel work in the unified
bootstrap and the paired mean-specification comparison. The resample indices
are constructed before the parallel callbacks. Worker count is not part of
semantic cache identity and does not by itself change an artifact path or the
intended scientific design.\\
Bootstrap execution mode: \code{HETID_BOOT_MODE} & Unset or empty selects
\code{reuse}. The only values are \code{reuse} and \code{rerun}, compared after
trimming and conversion to lower case. & \code{reuse} accepts only one
semantically valid, all-or-nothing unified cache; a missing or invalid cache
warns and causes a complete rerun. \code{rerun} skips the cache read and
computes the complete stage. A newly computed candidate is validated and
transactionally replaces the cache before downstream publication. The mode
value itself is not semantic provenance.\\
Draft-run acknowledgement: \code{HETID_ALLOW_DRAFT_RUN} & Unset or empty is
the default. Any nonempty value acknowledges a replication count other than
10,000. & Releases only the pre-directory publishing guard. It does not select
the replication count, change an estimator or inference rule, or enter cache
identity. It permits a reduced-count run to overwrite tracked publication
files.\\
FRED transport selector: \code{FRED_API_KEY} & Unset or empty selects the
public CSV endpoint through the curl fallback. Any nonempty key selects the
official JSON API; the key is never logged. & Changes the acquisition route for
FRED series such as PCECC96. The API route and CSV route have separate checked-in
retry and timeout policies. The key is not cache provenance. If the acquired
data differ, the bootstrap input hash invalidates reuse; identical acquired
data leave the scientific input unchanged.\\
Quadratic equivalence assertion: \code{HETID_ASSERT_EQUIV} & Unset or empty
disables it; any nonempty value enables it. & Each reached generalized
quadratic-system build also runs the legacy builder and requires exact identity
of the solver-consumable quadratic list and its $L_i$, $V_i$, and $Q_i$ leaves.
A mismatch stops execution; a pass leaves scientific objects unchanged. The
flag is not cache identity. A cache hit skips bootstrap draw callbacks, so
\code{HETID_BOOT_MODE=rerun} is needed when the intended check includes those
calls.\\
Tracked-artifact retention during reset: \code{--keep-tracked} & The flag is
supplied to \file{scripts-paper/reset_pipeline_state.R} and is absent by default.
The reset entrypoint is separate from the pipeline runner. & Every reset clears
the unified-bootstrap cache, gate state, diagnostics, and LaTeX sidecars. Without
the flag it also clears tracked tables, figures, and reports. With the flag,
those three publication groups are restricted to gitignored paths, so tracked
publication files survive. The flag changes only what state a later run
inherits.\\
Mean-specification plan: \code{PAPER_SPEC_PLAN} &
The tracked plan in \file{scripts-paper/config/analysis.R} is
\code{mean=c("B","A")} and
\code{volatility="B"}. Each panel must contain one or more distinct A/B entries;
the first mean entry is published, and the volatility entry must equal it. &
Controls which mean fits are computed and which fit supplies every downstream
log-variance stage. A one-entry mean plan skips alternative-specification draws
but still writes the required published-only comparison CSV. Changing the first
entry changes the published system, all downstream estimators and inference,
and unified-bootstrap identity.\\
Joint-GMM decision: \file{config/decisions/joint_gmm.R} & The committed record
uses the no-answer default: both nondefault moment blocks and stage C are off,
the intercept target is \code{a_L}, and the moment-tolerance grid is empty. &
The runner performs graph, algebra, and software replication checks and writes
the required joint-GMM diagnostic pair. The schema accepts other values, but
the current driver stops with \code{unwired_joint_branch} on every otherwise
valid nondefault record. The decision therefore cannot enable substantive
joint-GMM work in this runner and does not affect unified-bootstrap identity.\\
EGARCH-X decision: \file{config/decisions/egarch.R} & The committed record binds
a fresh \code{non_reject} residual-dynamics verdict and sets both decisions to
\code{not_asked}. The schema also represents the rejecting and unreliable gate
branches listed in Table~\ref{tab:outside-runner}. & The runner requires exact
agreement between this record and the fresh residual-dynamics gate; disagreement
stops. Agreement with the committed record closes the dynamic workstream and
requires all four conditional EGARCH-X artifacts to be absent. No record can
enable EGARCH-X estimation in the current runner because it injects no dependency
state and calls no estimator or artifact producer. The record does not affect
unified-bootstrap identity.\\
LAD dependency decision: \file{lad.dcf} & The committed record at
\file{scripts-paper/config/decisions/lad.dcf} has decision
\code{approved} for \code{quantreg} 6.1. Decision values are
\code{unanswered}, \code{declined}, and \code{approved}; consent values are
\code{not_asked}, \code{not_needed}, \code{declined}, and \code{approved}. A
missing path passed directly to the reader is classified as unanswered, but the
runner's required \code{paper_path()} lookup stops if the tracked file is absent. &
An existing unanswered or declined record prevents LAD source loading. Approval requires the installed package version to equal the
recorded version; absence or mismatch stops execution. The choice governs the
three conditional LAD artifacts. LAD is outside the unified bootstrap, so the
choice does not invalidate that cache.\\
\bottomrule
\end{longtable}
\endgroup

The six environment-variable rows are exhaustive: no other production path
rooted at \file{scripts-paper/run_pipeline.R} reads a user environment variable.
The four tracked rows exhaust the top-level plan and decision records consumed
by the runner. The reset flag is the only command-line option of the reset entrypoint. The
validation table comparator has its own command-line interface, but it is a
separate validation utility and cannot change pipeline execution.

The source also contains global constants, protocol vocabularies, and nested control objects. A user
can edit those files, but the runner does not expose their leaves as supported
runtime options. This document calls them \emph{source-fixed contracts}. The
distinction matters because such an edit changes the code being run and may
invalidate the cache at file or whole-object granularity. Table~\ref{tab:source-controls}
accounts for every source-fixed control family in the production graph.

\begin{landscape}
\begingroup
\footnotesize
\begin{longtable}{@{}>{\raggedright\arraybackslash}p{0.28\linewidth}
  >{\raggedright\arraybackslash}p{0.27\linewidth}
  >{\raggedright\arraybackslash}p{0.37\linewidth}@{}}
\caption{Complete map of source-fixed global control and protocol families. Paths are
relative to \file{scripts-paper/}. The detailed sections state current
scientific values; this table states what a source edit would affect.
\label{tab:source-controls}}\\
\toprule
Authority family & What it controls & Downstream and cache consequence\\
\midrule
\endfirsthead
\toprule
Authority family & What it controls & Downstream and cache consequence\\
\midrule
\endhead
Loose settings in \file{config/analysis.R} & \code{fred_from}, \code{fred_to},
\code{date_begin}, \code{date_end}, \code{show_mats}, \code{lag_qtrs}, and
\code{boot_seed} fix the pull dates, analysis dates, displayed maturities, lag,
and bootstrap seed. \code{news_prefix} and \code{expected_prefix} fix the
generated-series names.
The mean-specification plan is the supported control listed separately in
Table~\ref{tab:user-controls}. & Changes acquisition span, sample rows, time
alignment, generated names, reported variables, or random resampling. Changes
that alter the stage input, design, or resampling invalidate unified-bootstrap
reuse. A prefix-only edit renames generated columns used by reports; the
prefixes are outside the draw-code and semantic-runtime hashes, so that edit
alone does not invalidate reuse.\\
\src{config/analysis_contract.R}{PAPER_ANALYSIS_CONTRACT} and
\code{PAPER_YIELD_VOL_SUFFIX} & PC dimensions and
preprocessing; relaxation-value axes; variance-share, inference, and figure
policies; ACM, yield-volatility, consumption, and instrument inputs; and the
generated yield-volatility suffix. & Propagates
through data construction, identified sets, both log-variance estimators,
bootstrap inference, diagnostics, and publication. The whole contract enters
the runtime hash, and its file is in the draw-code manifest.\\
Path, artifact, lifecycle, and publication contracts in \file{config/paths.R},
\file{config/artifacts.R}, the \file{config/artifact_manifest_*.R} builders,
\file{config/artifact_lifecycle.R}, \file{config/artifact_latex.R}, and
\file{config/artifact_reset.R} & The output root;
artifact groups, IDs, paths, lifecycle statuses, owners, lookup families and
variants; generated sweep and region records; LaTeX labels and publication
triples; conditional-route schema; and reset compartments. & Change where files
are written, which files are expected, how they are addressed and cleaned, and
which TeX/PDF triples must compile. The runner always loads and validates these
contracts before creating directories. Artifact lifecycle status does not select
execution reachability. These auxiliary artifact files are not themselves all
members of the bootstrap draw-code manifest, so an artifact-only source edit need
not invalidate existing result draws; a changed cache path can instead move where
the one cache artifact is read or written.\\
\src{config/acquisition.R}{PAPER_FRED_DOWNLOAD_CONTROL} & FRED endpoints,
attempts, backoff, curl flags, retries, and timeout/stall rules. & Changes only
data transport and failure behavior directly. Any resulting data change reaches
all downstream estimation and the input hash. The API-key environment variable
is the supported transport selector listed separately above.\\
\src{config/diagnostics.R}{PAPER_HETEROSKEDASTICITY_CONTROL} & Diagnostic test
suites, Goldfeld--Quandt settings, rejection threshold, and caption tests. &
Changes diagnostic computation and its published table, not estimator choice
or unified-bootstrap draws.\\
\src{config/figure_rendering.R}{PAPER_FIGURE_STYLE} & Colors, line widths,
point sizes, opacity, markers, and region fills. & Changes rendered bytes and
appearance, not estimates or cache validity.\\
\src{config/figure_rendering.R}{PAPER_FIGURE_RENDER_CONTROL} & Projection and
three-dimensional grids, views, region relaxation values, device sizes, labels,
fonts, and layouts. & Changes geometry evaluation and figure content. Region
relaxation values also change manifest membership; the fitted-volatility sweep
axis instead belongs to the analysis contract. This file is outside the
unified-bootstrap draw-code hash.\\
\src{config/inference_search_control.R}{PAPER_INFERENCE_SEARCH_CONTROL} &
Imbens--Manski and Stoye calibration roots, bivariate-normal integration,
robust endpoint-correlation diagnostics, boundedness-frontier search, and
bootstrap failure and progress rules. & Can change confidence-interval
calibration, endpoint-correlation estimates, frontier classification,
bootstrap failure gates, and progress messages. The whole object enters the
runtime hash and its file enters the draw-code hash, so any edit invalidates
reuse.\\
\src{support/statistics/mbb_protocol_authority.R}{paper_mbb_protocol} & The
circular-MBB resampler identifier, primary block-length rule
$\ell=\lceil1.5T^{1/3}\rceil$, RNG kinds, resample-index family identities, and
family schema. & Supplies the primary block length from which
\code{bootstrap_stage_sensitivity_block_length()} derives the doubled
volatility-sensitivity length. The index-family code implements circular
resampling for both result branches and the paired mean-specification
comparison. This protocol file is draw-hashed; the design, resampler, block
rule, RNG, and index and post-index RNG hashes also enter semantic provenance,
so a change invalidates reuse.\\
\src{config/inference_search_control.R}{PAPER_LOGVAR_BUDGETS} & The identical
primary PPML/Harvey set-map budgets; the fresh 40,000-fit budget used by each
second search audit; PPML's larger 8,000-row selected-grid cap for that audit;
and the per-draw and fitted-volatility budgets. & Changes how much numerical
search each named task may perform and whether it stops at a budget. PPML and
Harvey have the same primary budget and the same second-search-audit fit budget;
only PPML's second audit expands the selected-grid cap. Because the whole
authority file is in the draw-code manifest, editing any budget invalidates
reuse, including a budget used only outside the unified stage.\\
PPML second-search-audit traversal policy
\src{log_variance/estimators/ppml/pilot_and_grid.R}
{LOGVAR_PPML_COVERAGE_PROTOCOL} & The protocol fixes schema and selector
identity and tells the engine to traverse the selected rows in their supplied
order. \code{logvar_ppml_morton_select()} in the same file chooses those rows;
the selected-grid cap is in \code{PAPER_LOGVAR_BUDGETS}. & Changes only PPML's second search audit;
it does not change the common PPML/Harvey set-map budget or create Harvey-only
scientific treatment. The authority lies inside the draw-hashed PPML directory,
so editing it still invalidates unified-bootstrap reuse.\\
\src{config/reporting.R}{PAPER_REPORTING_CONTROL},
\code{PAPER_TABLE_STYLE}, \code{PAPER_LATEX_CONTROL}, and the shared missing-cell
and notes-label tokens & Covariance choices,
significance thresholds, displayed precision, missing-value tokens, table
layout, compilation, and sidecar cleanup. & Changes reported uncertainty,
stars, text, layout, or LaTeX lifecycle. Because \file{config/reporting.R} is in
the draw-code manifest and the reporting object enters the runtime hash, even a
presentation-only edit in this file currently invalidates reuse.\\
\src{config/statistics.R}{PAPER_STATIONARITY_CONTROL} & Summary-ACF maximum,
ADF deterministic term and lag selection, response-surface settings, KPSS type,
and Ljung--Box lag and type. & Has helper
consumers but no current production caller, so it changes no artifact in the
present runner unless that path becomes reachable.\\
\src{support/identification/quadratic_evaluation.R}{PAPER_QUADRATIC_CONTROL} &
Quadratic symmetry and feasibility tolerances, solver boxes and limits, point
identification, boundedness, and admission rules. & Can change feasibility,
endpoint status, and every dependent identified set or set map. The entire
identification directory is draw code, so edits invalidate reuse.\\
Log-variance estimator and search authorities in
\file{config/logvar_estimators.R} and
\file{log_variance/estimators/controls.R} &
\code{PAPER_LOGVAR_ESTIMATORS}, \code{LOGVAR_SEARCH_CONTROL}, and the
estimator-specific \code{LOGVAR_*_CONTROL} objects fix response scales,
capabilities, optimizer settings, multistart/search rules, and LAD
search/refinement. & Change which
estimators exist, how their sets are searched, numerical statuses, and which
tables and figures can be produced. Primary PPML and Harvey code and controls
are draw code. LAD computation remains outside the unified bootstrap, but its
controls share the draw-hashed \file{controls.R}; editing them therefore
invalidates reuse at file granularity. The offline LAD refinement control is
defined in that file but has no current production caller.\\
Residual-dynamics gate policy
\src{log_variance/diagnostics/protocols.R}{LOGVAR_DYNAMICS_GATE_PROTOCOL} and
the gate schema in \file{log_variance/diagnostics/dynamics/gate_core.R} &
Tested lags, the decision lag, significance level, maximum reported ACF lag,
and tie tolerance. & Changes the
diagnostic verdict and EGARCH-X routing evidence, not the currently published
estimator or unified-bootstrap cache.\\
Joint-null diagnostic controls
\src{log_variance/diagnostics/protocols.R}{LOGVAR_JOINT_NULL_CONTROL} and the
schemas in \file{log_variance/diagnostics/joint_null/schema.R} &
Search starts, tolerances,
stability perturbations, and certificate rules for the joint-zero restriction. &
Change whether a
compatible joint-null witness is found and how it is certified; they do not
add an identifying restriction to the main estimator. This diagnostic lies
outside the unified bootstrap and does not affect its cache.\\
Joint-GMM numerical controls and budgets
\src{log_variance/diagnostics/joint_gmm/epigraph.R}
{logvar_joint_gmm_constants}, together with the row and artifact schemas & The current default graph replication directly
uses its replication-grid size, point cap, root tolerance, and design-rank
tolerance. The complete constants object is stamped into the diagnostic
specification identity; remaining optimizer and search settings belong to
nondefault branches. & The grid size and point cap change replication points,
the root tolerance changes its attained/unreliable verdict, and the design-rank
tolerance changes the current fail-closed design check. Any constant edit
changes the diagnostic specification identity and record, but unwired-branch
settings do not change current arithmetic.
Joint-GMM lies outside the unified bootstrap and does not affect its cache.\\
Joint-GMM routing decision
\src{log_variance/diagnostics/joint_gmm/decision_schema.R}
{LOGVAR_JOINT_DECISION_SCHEMA_VERSION} and
\src{config/decisions/joint_gmm.R}{logvar_joint_gmm_decision} & The schema,
decision-specification identity, enabled moment blocks, intercept target,
stage-C request, moment-tolerance grid, Gaussian-gap restriction, provenance,
and rationale. & The driver validates the checked-in record before computing
the diagnostic. The current record selects the no-answer default: both nondefault
moment blocks and stage C are off, the intercept target is \code{a_L}, and the
tolerance grid is empty. Any otherwise valid nondefault record stops with
\code{unwired_joint_branch}; it does not silently change current arithmetic.
The full record enters the joint-GMM RDS; its decision-specification identity
enters the flat status rows. This diagnostic lies outside the unified bootstrap
and does not affect its cache.\\
EGARCH-X routing decision and protocol
\src{log_variance/extensions/egarch/decision_core.R}{LOGVAR_EGARCH_*} and
\src{config/decisions/egarch.R}{logvar_egarch_decision} & The decision schema,
fresh-gate binding, plan and upstream-plan hashes, approved dependency version,
approval-prompt hashes, closed decision/provenance/terminal-status values, and
the checked-in routing record. & Before routing, the validator binds the record
to the fresh residual-dynamics gate and hard-stops on a stale hash, schema,
prompt, or inconsistent decision ladder. A valid record determines the route's
branch together with dependency availability and version. The status record
stores selected decision fields, the resulting terminal status, and the
\code{run_dynamic} value. The current non-reject record leaves both approvals
\code{not_asked}; the current runner has no dynamic producer. This route lies
outside the unified bootstrap and does not affect its cache.\\
\src{support/runtime/core.R}{PAPER_SERIALIZATION_CONTROL}, consumed by
\file{support/artifacts/typed_artifacts.R} & RDS version, numeric
serialization, CSV typing, readback tolerance, escaping, encoding, and
year-quarter formatting and validation. & Change
artifact encoding and validation. These files participate in stage
specification or draw-code identity, so relevant edits invalidate reuse.\\
Unified-bootstrap cache and provenance contracts in
\file{inference/bootstrap_stage_specs.R},
\file{inference/bootstrap_stage_cache.R},
\file{support/inference/bootstrap_stage_code_manifest.R}, and
\file{support/statistics/boot_freshness.R} & Stage-specification and cache
schemas and required fields;
legacy cache names; draw-code and presentation source manifests; SHA-256 field
rules; and the R, platform, BLAS, LAPACK, and selected-package runtime
fingerprint. & Determine whether a payload is structurally and semantically
eligible for reuse and which older cache paths are removed after validation.
Changing the draw-code manifest can change \code{code_sha}; changing only the
presentation manifest changes its audit hash without invalidating result draws.
These contracts choose cache acceptance, not an estimator.\\
Unified-bootstrap complete-case policy
\src{inference/bootstrap_stage_specs.R}{BOOTSTRAP_STAGE_COMPLETE_CASE_POLICY} &
Records and validates the current timing and roles for applying
\code{stats::complete.cases} to nominal asset-return-PC columns after shared
mean estimation. The filtering helper currently hard-codes that predicate and
uses the policy's role names as an assertion. & The current implementation
keeps every resampled row in the shared mean branch, then subsets the
log-variance branch's \code{w1}, \code{w2}, quarter key, and PC data to rows
with complete PCs. Editing the policy object alone changes stage identity and
invalidates reuse, and may fail its role-name assertion; it does not change
selected rows unless the filtering helper changes too.\\
Closed status and normalization contracts in
\file{support/identification/status_contract.R},
\file{log_variance/engine/contracts.R}, and
\file{support/statistics/normalizations.R} & The allowed endpoint and fit-status
vocabularies, engine phases, and Gaussian log-square or median-scale constants. &
These are protocol definitions and mathematical normalizations, not supported
user options. A source edit changes which states the engine can accept or how
scale/intercept conversions are normalized; the files are draw-hashed, so the
edit invalidates unified-bootstrap reuse.\\
\bottomrule
\end{longtable}
\endgroup
\end{landscape}

The tables establish the change path. Later sections explain the affected
scientific objects and artifacts in runner order. The cache rule is not simply
scientific versus cosmetic: the implementation hashes named files and
whole control objects. Consequently, a source edit in a draw-hashed file can
invalidate reuse even when the edited field affects only presentation.

\subsection{Functionality outside the current runner}
\label{sec:outside-runner}

Table~\ref{tab:outside-runner} prevents a source file, formal argument, artifact
record, or command from being mistaken for a branch that
\file{run_pipeline.R} executes.

\begingroup
\footnotesize
\begin{longtable}{@{}>{\raggedright\arraybackslash}p{0.25\textwidth}
  >{\raggedright\arraybackslash}p{0.23\textwidth}
  >{\raggedright\arraybackslash}p{0.44\textwidth}@{}}
\caption{Alternative and separate functionality. Status uses the exact terms in
Section~\ref{sec:status-words}; mixed implementations receive separate rows.
\label{tab:outside-runner}}\\
\toprule
Functionality and source & Status & Exact boundary and way to invoke or enable it\\
\midrule
\endfirsthead
\toprule
Functionality and source & Status & Exact boundary and way to invoke or enable it\\
\midrule
\endhead
ACM data-contract choices not selected by the paper contracts &
Alternative path, source-selected &
The active runner calls \code{hetid::extract_acm_data()} for both quarterly and
daily data. A source edit can make those calls request GitHub or New York Fed
monthly data explicitly, request \code{risk_neutral_yields}, or drop incomplete
quarters. The checked-in contracts instead request quarterly yields and term
premia with incomplete quarters kept and daily yields, both through
\code{source="auto"}. No runner environment input or command-line flag selects
the alternative path.\\
Direct ACM downloader package API &
Separate entrypoint &
\code{hetid::download_term_premia()} exposes source, frequency, force, and quiet
arguments to a direct package caller. This direct call is separate from the paper
runner. Within the selected daily \code{source="auto"},
\code{auto_download=TRUE} path, however, the active extractor can call the same
download implementation at runtime when no valid daily asset is available.\\
Nondefault joint-GMM moment, profile, and search branches &
Available but not called &
The schema and source-loaded modules represent the $z$ and log/PPML moment
blocks, intercept target \code{a_P} or both targets, stage C, and a nonempty
moment-tolerance grid. The production driver validates the record and then stops
with \code{unwired_joint_branch} before computing any of these branches. The
tracked record is a top-level behavior input, but none of these fields can enable
the source-loaded substantive code in the current runner.\\
Joint-GMM Gaussian-gap restriction &
Not implemented &
The decision schema reserves the flag but rejects \code{TRUE} with
\code{unsupported_gaussian_gap_restriction}. No Gaussian-gap solver exists.\\
EGARCH-X terminal stop routes &
Alternative path, source-selected and runtime-selected outcome &
For a fresh unreliable gate, a matching record produces \code{gate_unreliable}.
For a fresh rejecting gate, matching decision records can produce
\code{no_answer}, \code{estimand_declined}, \code{dependency_declined}, or
\code{reject_awaiting_answer}. The committed record instead binds
\code{non_reject} and produces \code{gate_non_reject} only if the fresh gate
matches it. A mismatch stops before routing.\\
EGARCH-X approved router result &
Available but not called &
The pure router can return \code{approved_pending_dynamic} after a rejecting gate,
two approvals, and an injected exact dependency version. The current runner
defines neither dependency-injection object, so its driver supplies unavailable
and missing-version defaults; no plain or source-decision path reaches this
result.\\
EGARCH-X estimator and artifact producer &
Not implemented &
The extension contains cleanup, decision, and routing code only. It contains no
EGARCH-X estimator or artifact producer, and the runner has no dynamic source line
after the router. The four conditional artifact records are reserved placeholders;
the current graph deletes their paths and requires them to remain absent.\\
Stationarity and general reporting helpers in
\file{support/statistics/bootstrap_and_stationarity.R} and
\file{support/statistics/reporting_and_validation.R} &
Available but not called &
The configuration loads definitions for summary statistics, ADF, KPSS, and
Ljung--Box stationarity reports and formatted or interactive tables. The current
production graph has no caller. Packages used only inside those uncalled
functions, including \code{moments} and \code{urca}, are not prerequisites for
the current runner merely because their qualified names appear in a loaded file;
the same is true of \code{cli} in the uncalled reporting helpers. Only these
named helpers are uncalled: the same \file{bootstrap_and_stationarity.R} file also
defines \code{mbb_index()}, which is production-active as the circular-MBB
resampler.\\
Source-loaded convenience APIs &
Available but not called &
Examples are artifact basename/directory accessors, \code{fine_tau_grid}, a
complete-case-policy assertion wrapper, an EGARCH cleanup compatibility helper,
and PPML/Harvey standard-error frame wrappers. Their owning production files are
loaded, but no production call site invokes these functions. They remain test or
future APIs, not hidden runner options.\\
PPML and Harvey mathematical convenience primitives &
Available but not called &
\code{logvar_ppml_score}, \code{logvar_ppml_info}, \code{logvar_harvey_objective},
\code{logvar_harvey_score}, and \code{logvar_harvey_moment} expose formulas for
tests or source-loaded future moment branches. The active PPML acceptance code
computes its score and information matrix directly; the active Harvey solver
uses \code{hv_eval}.
Section~\ref{sec:logvar} maps the equations to those active callers.\\
Legacy log-OLS caption
\file{log_variance/tables/legacy_log_ols_caption.R} &
Available but not called &
Retains a caption builder for a separate legacy mean-log OLS robustness table.
The runner never sources this file.\\
LAD offline refinement
\file{log_variance/estimators/lad/offline_refinement.R} &
Available but not called &
Only the LAD outer-map test suite sources it. The selected production LAD path
does not.\\
Bounds-figure test support
\file{log_variance/figures/bounds_by_tau_test_support.R} &
Available but not called &
Supplies plot-data support to tests. The production bounds renderer never sources
it. These three inactive files are the production-tree exclusions listed in
\file{scripts-paper/README.md}.\\
Full state reset \file{scripts-paper/reset_pipeline_state.R} &
Separate entrypoint &
\code{Rscript scripts-paper/reset_pipeline_state.R} clears the cache, gate state,
diagnostics, publication outputs, and sidecars by compartment. Its
\code{--keep-tracked} flag retains tracked publication files while still clearing
eligible untracked state. The pipeline never calls this command.\\
Published-table comparator
\file{scripts-paper/validation/compare_output_tables.R} &
Separate entrypoint &
It requires exactly two positional output-root arguments, projects published
numeric TeX table content, and reports equality or differences. It does not run
the pipeline, select a branch, or validate non-table artifacts.\\
Paper test orchestrator \file{scripts-paper/tests/run_tests.R} &
Separate entrypoint &
It launches two structural checks and 36 listed test suites in fresh R
processes. The runner does not call the orchestrator or any test entrypoint.
Individual test fixtures may select helper arguments that production callers
fix; those fixture choices are not pipeline options.\\
Package quality suite \file{docs/quality-check.R} &
Separate entrypoint &
\code{Rscript docs/quality-check.R} runs package-quality tools and writes quality
reports under \file{docs/}. It belongs to neither the paper production graph nor
the paper test orchestrator.\\
\bottomrule
\end{longtable}
\endgroup

\subsection{Publishing-run guard}

\src{scripts-paper/config/analysis.R}{resolve_whole_number_env} parses
\code{HETID_BOOT_REPS} and \code{HETID_BOOT_CORES} as whole numbers. Subsequent
configuration assertions reject values below 2 and 1, respectively. The default
replication count is 10,000.
The platform-dependent worker default is stated in Table~\ref{tab:user-controls}.
Before the runner creates directories or deletes conditional artifacts, it
rejects a different count unless \code{HETID_ALLOW_DRAFT_RUN} is nonempty.
\code{HETID_ALLOW_DRAFT_RUN} acknowledges that a reduced run rewrites
publication files; it neither selects the count nor changes the scientific
publication standard of 10,000.

\src{scripts-paper/config/analysis.R}{resolve_boot_mode_env} maps
\code{HETID_BOOT_MODE} to \code{PAPER_BOOT_MODE}. The default is \code{reuse};
the only alternative is \code{rerun}. This choice is consulted later by the
unified-bootstrap stage. It does not weaken the publishing-run guard.

A reduced serial run needs a reduced replication count, the draft
acknowledgement, and an explicit one-core request:
\begin{verbatim}
HETID_BOOT_REPS=8 HETID_BOOT_CORES=1 HETID_ALLOW_DRAFT_RUN=1 \
  Rscript scripts-paper/run_pipeline.R
\end{verbatim}
The acknowledgement and replication count are the settings required by the
publication guard; \code{HETID_BOOT_CORES=1} separately requests serial
execution. A production serial command is:
\begin{verbatim}
HETID_BOOT_REPS=10000 HETID_BOOT_CORES=1 \
  Rscript scripts-paper/run_pipeline.R
\end{verbatim}
Leaving \code{HETID_BOOT_REPS} unset also selects 10,000 replications. The
commands are interpretive examples; this review executed no pipeline command.

\subsection{Artifact manifest directories and startup cleanup of conditional artifacts}

After the guard, \src{scripts-paper/config/artifacts.R}{create_artifact_directories}
creates each unique parent in the expanded artifact manifest. The runner then calls
\src{scripts-paper/config/artifact_lifecycle.R}{cleanup_conditional_artifacts}
for \code{conditional_lad} and \code{conditional_egarch}. The cleanup helper
refuses \code{required}, removes every path whose artifact lifecycle status is
the selected conditional status, rereads path existence, and fails if any target
survives.

The runner preserves required artifacts until their producers overwrite them,
so an early failure may leave files from unreached stages of an earlier run.
The separate \file{scripts-paper/reset_pipeline_state.R} entrypoint owns
full-reset behavior.

Invoking that reset entrypoint without a flag removes the unified-bootstrap
cache, gate state, diagnostics, tables, figures, reports, and LaTeX sidecars.
Adding \code{--keep-tracked} still removes the cache, state, diagnostics,
sidecars, and gitignored publication artifacts, but retains tracked table,
figure, and report files. The runner itself never calls this full reset.

\section{Data construction}
\label{sec:data}

The pipeline builds quarterly analysis series for 1962 Q1--2025 Q4. It combines
validated quarterly ACM yields and term premia, FRED PCECC96 consumption, daily
ACM yields, the bundled \code{hetid::variables} nominal asset-return PC scores,
and the SDF-news, expected-SDF, and lagged expected-SDF panels. The quarterly ACM contract sets
\code{frequency="quarterly"}, \code{auto_download=FALSE}, \code{source="auto"},
and \code{use_incomplete_quarters=TRUE}. The daily ACM yield input uses
\code{frequency="daily"}, \code{auto_download=TRUE}, and \code{source="auto"};
its cache-only asset downloads on first use. Daily ACM yields and their realized
quarterly volatility use percentage points. FRED PCECC96 growth is a
nonannualized quarterly percent rate. Final sample size is a runtime fact outside
this source-only trace.

The data-construction source modules appear in their literal runner order.

\subsection{SDF-news, expected-SDF, and lagged expected-SDF panels}

\src{scripts-paper/data_preparation/build_sdf_series.R}{sdf_news} consumes the
validated quarterly ACM matrices. For maturities of 3--117 months, it calls the
package SDF-news and expected-SDF kernels. The analysis stage produces
\code{sdf_news}, \code{expected_sdf}, and \code{lag_expected_sdf}; the last
object relabels expectations so the previous-quarter expectation aligns with
the next realized outcome.

Quarterly yield and term-premium levels enter through this SDF construction.
Daily yields separately construct the five-year-yield volatility instrument in
the daily-yield analysis stage. The bundled \code{hetid::variables} data supply the nominal
asset-return PCs.

\subsection{Quarterly real consumption growth}

\src{scripts-paper/data_preparation/build_consumption_growth.R}{gr1_pcecc96}
retrieves FRED series \code{PCECC96}, normalizes dates to period end, and computes
\[
  100\left(\frac{C_t}{C_{t-1}}-1\right).
\]
Here $C_t$ is FRED PCECC96 in quarter $t$. The formula reports percent growth
and is not annualized.
The configured FRED pull begins in 1947 to provide needed lags; downstream
analysis uses the 1962 Q1--2025 Q4 window.

\code{FRED_API_KEY} selects the acquisition route. A nonempty key uses the
official FRED JSON API; an empty or missing key uses the public CSV endpoint
through curl with the checked-in HTTP/1.1, retry, and stall-timeout policy. The
key is never logged. The route itself is not bootstrap provenance, but any
difference in the acquired PCECC96 observations changes the stage input hash
and invalidates unified-bootstrap reuse.

\subsection{Yield-volatility panel and instrument}

\src{scripts-paper/data_preparation/build_yield_volatility.R}{yield_vol}
loads daily ACM yields, sums squared daily yield changes within each quarter,
and defines the square root as the yield-volatility panel. Its five-year column
\code{y60_vol} supplies the five-year-yield volatility instrument for all three
SDF-news PC scores. The identification estimator demeans the instrument on its
aligned sample.

\subsection{Nominal asset-return PCs}

\src{scripts-paper/data_preparation/build_asset_return_pcs.R}{lag_asset_return_pc}
reads the bundled nominal-financial-asset-return PCs \code{pc1}--\code{pc4},
normalizes their period-end dates, and relabels them one quarter forward as
\code{l.pc1}--\code{l.pc4}. The stage uses the bundled PCs directly; yield or
term-structure PCs do not exist in this pipeline. Later log-variance analysis
stages center the four scores without scaling.

\subsection{Three SDF-PC score blocks}

\src{scripts-paper/data_preparation/build_sdf_pcs.R}{pc_scores} calls
\src{scripts-paper/config/analysis_contract.R}{paper_model_pca} with
\code{center=TRUE} and \code{scale.=TRUE} separately for the expected-SDF,
lagged expected-SDF, and SDF-news panels after windowing and complete-case
filtering. It constructs the three SDF-PC score blocks. It aligns the sign of
each lagged expected-SDF PC score with the corresponding shifted expected-SDF
PC score. The products are
\code{expected_sdf_pc}, \code{lag_expected_sdf_pc}, and \code{sdf_news_pc}.
These stages leave the aligned outcome, five-year-yield volatility instrument,
and three SDF-PC score blocks consumed by the mean-equation stages.

\section{Structural mean equation for quarterly real consumption growth and set identification}
\label{sec:mean}

Aligned data enter two mean specifications and one shared moment and quadratic
system. Published mean specification B residualizes each SDF-news PC score on
$X$. Diagnostic mean specification A imposes a zero reduced form.

\subsection{Exogenous-news OLS benchmark and identification sample}

\src{scripts-paper/mean_equation/fit_ols.R}{ols_mean_eq} full-joins quarterly
real consumption growth, the lagged expected-SDF PC score block, and the
SDF-news PC score block by quarter. It fits
the exogenous-news benchmark
\[
  \Delta c_t = X_t\beta_1^{\mathrm{OLS}}
    + Y_{2t}\theta^{\mathrm{OLS}} + u_t.
\]
Here $\Delta c_t$ is quarterly real consumption growth. $X_t$ contains an
intercept and the lagged expected-SDF PC score block, and $Y_{2t}$ contains the
SDF-news PC score block. The vectors $\beta_1^{\mathrm{OLS}}$ and
$\theta^{\mathrm{OLS}}$ are the corresponding OLS coefficients; $u_t$ is the
exogenous-news OLS benchmark residual.

The identified-set driver builds its own frame. It adds \code{y60_vol}, applies
the analysis window, drops incomplete rows, sorts by quarter, and retains this
aligned data inside each fit. The OLS and identification objects own their row
policies and do not rely on ambient row positions.

\subsection{Published mean specification B and diagnostic mean specification A}
\label{sec:specs}

\src{scripts-paper/config/analysis.R}{PAPER_SPEC_PLAN} sets the mean plan to
\code{c("B", "A")}, with published mean specification B first and diagnostic
mean specification A second. The driver verifies that the log-variance plan
contains only B, retains both fits in \code{set_id_mean_eq_by_spec}, and binds B
to the downstream name \code{set_id_mean_eq}.

This tracked plan is the supported source-level specification control. Each
panel must contain at least one distinct entry drawn from A and B. The first
mean entry is published, and the sole volatility entry must equal it. A
one-entry mean plan skips the alternative-specification draws but still writes
the required comparison CSV with published-specification rows. Changing the
first entry changes the published mean system, every downstream log-variance
fit and inference object, and unified-bootstrap identity.

Both specifications first regress quarterly real consumption growth on $X$ and
call the residual $W_1$. Their only structural difference is the reduced form
for $Y_2$.
\begin{table}[htbp]
\centering
\small
\caption{Published and diagnostic mean specifications. SDF means stochastic
discount factor, and PC means principal component. The mean-control block $X$
contains an intercept and the lagged expected-SDF PC scores. $Y_2$ is the
SDF-news PC score block. $W_2$ is the specification-specific news block after
the stated reduced-form operation. $\beta_2^R$ is the news-on-controls
coefficient matrix. $\beta_1(\theta)$ is the control coefficient at mean news
coefficient $\theta$. Both specifications use the same outcome, controls,
instrument, seven centered identification moments, and quadratic-system
machinery. They differ only in the reduced form for $Y_2$.
\label{tab:mean-specs}}
\begin{tabular}{@{}p{0.16\textwidth}p{0.35\textwidth}p{0.35\textwidth}@{}}
\toprule
 & Published mean specification B & Diagnostic mean specification A\\
\midrule
\code{impose_null} & \code{FALSE} & \code{TRUE}\\
$W_2$ & residuals from regressing each SDF-news PC score on $X$ & raw SDF-news
PC score block $Y_2$\\
$\beta_2^R$ & estimated news-on-controls coefficient matrix & exact zero matrix\\
$\beta_1(\theta)$ &
$\beta_1^R-(\beta_2^R)\trans\theta$ &
constant $\beta_1^R$\\
Use & all published mean and log-variance analysis stages & paired B-versus-A
mean-specification comparison\\
\bottomrule
\end{tabular}
\end{table}

Table~\ref{tab:mean-specs} shows the analysis roles directly. Published mean
specification B feeds all published mean and log-variance stages. Diagnostic
mean specification A enters only the paired comparison.
The shared implementation is
\src{scripts-paper/support/identification/identified_set_bootstrap.R}
{estimate_set_id_system}. It demeans $Z=\code{y60_vol}$, computes the package
moments, builds the tau-zero system, and solves the tau-zero point when the
system has the required rank and consistency.

\subsection{Triangular system}

After the estimator constructs $W_1$ and the specification-specific $W_2$, the
identification problem is
\begin{align*}
  W_{1t} &= W_{2t}\theta + \varepsilon_{1t},\\
  W_{2t} &=
  \begin{cases}
    \text{SDF-news reduced-form residuals}, & \text{published mean specification B},\\
    Y_{2t}, & \text{diagnostic mean specification A}.
  \end{cases}
\end{align*}
Here $W_1$ is quarterly real consumption growth residualized on $X$. $W_2$ is
the specification-specific three-score block in the display. $\theta$ is the
mean news coefficient. $\varepsilon_1$ is the structural residual.
The code applies the relaxed heteroskedasticity restriction to each SDF-news PC
score $i$:
\[
  \left|\operatorname{Corr}
    \left(Z_t,\varepsilon_{1t}W_{2it}\right)\right|
  \leq \tau_i
  \left|\operatorname{Corr}\left(Z_t,W_{2it}^{2}\right)\right| .
\]
The restriction compares the instrument's correlation with
$\varepsilon_{1t}W_{2it}$ with its correlation with $W_{2it}^2$. At
$\tau_i=0$, it imposes the Lewbel orthogonality condition exactly. Positive
values allow a scaled departure, with the score's heteroskedasticity correlation
setting the scale. This restriction is the source's identifying assumption.
Economic validity requires evidence beyond this static code trace. At zero, the
system yields the tau-zero point when its linear equations are well determined.
Positive common relaxation values
produce a joint set for the three-vector $\theta$. The paper sets
\code{gamma = matrix(1, 1, 3)}: the single demeaned instrument enters each of
the three score-specific restrictions with unit weight.

\subsection{Seven centered identification moments}
\label{sec:moments}

These seven centered moment blocks convert the relaxed correlation restrictions
into inputs to the quadratic identified-set system.
\src{R/hetid_moments.R}{compute_identification_moments} returns a validated
\code{hetid_moments} container. All covariances use
\src{R/statistics_utils.R}{centered_cov}:
\[
  \Covhat(A,B)=\frac{1}{T}(A-\bar A)\trans(B-\bar B).
\]
$A$ and $B$ are aligned sample vectors or matrices. $T$ is their common row
count. Bars denote column means. The divisor is $T$ rather than $T-1$. For $W_{2i}$, let
$W_2\odot W_{2i}$ denote the matrix whose columns multiply each $W_2$ column by
$W_{2i}$. The package computes
\begin{align*}
  S_i^{(0)} &= \Varhat(W_1W_{2i}),&
  \sigma_i^2 &= \Varhat(W_{2i}^2),\\
  R_i^{(0)} &= \Covhat(Z,W_1W_{2i}),&
  R_i^{(1)} &= \Covhat(Z,W_2\odot W_{2i}),\\
  P_i^{(0)} &= \Covhat(Z,W_{2i}^2),&
  S_i^{(1)} &= \Covhat(W_1W_{2i},W_2\odot W_{2i}),\\
  &&S_i^{(2)} &= \Varhat(W_2\odot W_{2i}).
\end{align*}
The container records the three-coordinate $\theta$ axis and the active
constraint indices. Downstream package functions validate these attributes;
alignment follows them exclusively.

\subsection{Quadratic identified-set system}
\label{sec:quadratic}

The quadratic system turns the seven moment blocks and the relaxation value
into one inequality per SDF-news PC score.
\src{scripts-paper/support/identification/quadratic_system.R}
{build_pipeline_quadratic_system} calls
\code{hetid::build_general_quadratic_system()}. With the paper's one instrument
combination per SDF-news PC score in matrix $\gamma$, this path is numerically
equivalent to the legacy builder on the $L_i$, $V_i$, and $Q_i$ quantities
consumed by the solver. For each SDF-news PC score $i$,
let $\gamma_{\cdot i}$ denote column $i$ of the weight matrix. Then
\begin{align*}
  L_i &= \gamma_{\cdot i}\trans R_i^{(0)},&
  V_i &= \gamma_{\cdot i}\trans
    \left(P_i^{(0)}P_i^{(0)\trans}\right)\gamma_{\cdot i},&
  Q_i &= \gamma_{\cdot i}\trans R_i^{(1)},\\
  d_i &= \frac{\tau_i^2V_i}{\sigma_i^2},&
  A_i &= Q_iQ_i\trans-d_iS_i^{(2)},\\
  b_i &= -2L_iQ_i+2d_iS_i^{(1)},&
  c_i &= L_i^2-d_iS_i^{(0)}.
\end{align*}
The identified set is
\[
  \Theta(\tau)=
  \left\{\theta\in\R^3:
  \theta\trans A_i\theta+b_i\trans\theta+c_i\leq0
  \text{ for every }i\right\}.
\]
\src{scripts-paper/support/identification/tau_star.R}{tau_quadratic_system}
broadcasts the paper's scalar $\tau$ over all three SDF-news PC scores. Constraint values
follow the \code{hin <= 0} convention: nonpositive means feasible.

\code{HETID_ASSERT_EQUIV} turns the numerical-equivalence statement into a live
assertion. When the variable is nonempty, every reached generalized build also
runs the legacy builder and requires exact identity of the quadratic list and
the $L_i$, $V_i$, and $Q_i$ leaves. A mismatch stops execution; a pass changes
no result. The flag does not enter cache identity. Therefore, a reused unified
cache does not recheck quadratic builds inside skipped bootstrap callbacks;
\code{HETID_BOOT_MODE=rerun} is required when those calls are part of the
intended assertion scope.

\subsection{Tau-zero point, profile endpoints, and the estimated boundedness frontier}

At $\tau=0$, \src{scripts-paper/support/identification/functional_bounds.R}
{solve_point_identification} solves the linear $Q\theta=L$ system. It requires
full column rank and consistency and records conditioning. At positive
relaxation values, coordinate and affine-functional endpoints use scaled
sequential least-squares quadratic programming (SLSQP) searches from
\file{scripts-paper/support/identification/profile_solver_core.R}. The profile
solver independently checks every returned point for feasibility and an active boundary. A
box-growth classifier separates a certified finite bound from an open direction.

\src{scripts-paper/support/identification/tau_star.R}{tau_star_fixed} locates the
bounded-to-unbounded transition. It brackets and bisects with certified
\code{bounded} and \code{unbounded} evaluations. An \code{unreliable} evaluation
does not count as evidence for either side.

The first profile solutions start at the origin. The current mean estimator then
uses \file{mean_equation/inference/theta_box_multistart.R} to add the tau-zero
point, axis, warm-chain, and cross-coordinate starts. It widens an endpoint only
with a certified endpoint-attaining point outside the current coordinate range.
The same endpoint-attaining points map through $\beta_1(\theta)$, so the reported
$\theta$ and $\beta_1$ boxes refer to the same attained set. The multistart
expansion certifies feasibility of added points; it
does not prove a global optimum for a nonconvex quadratic set.

\subsection{Relaxation-value grid and mean figures}

\src{scripts-paper/config/tau_grid.R}{paper_bounds_tau_grid} starts from the
historical 25-point backbone. It removes zero and the estimated boundedness
frontier and subdivides configured tail intervals. The resulting
relaxation-value grid is nonuniform and denser near the frontier.
\src{scripts-paper/mean_equation/inference/compute_bounds_by_tau.R}
{mean_eq_bounds_tau} stores full three-status theta tables under exact canonical
relaxation-value keys. It drops uncertified rows only from plotted ribbons,
writes \code{mean_bounds_figure}, and reinserts the refined displayed-relaxation-value
boxes used by the log-variance maps.

The projection renderer writes figures with and without the OLS benchmark.
The three-dimensional renderer crosses four relaxation values
$\{0.05,0.10,0.20,0.30\}$ with two unit systems and two OLS modes, producing
$4\times2\times2=16$ manifested SVGs.

\subsection{Attained variance-share ranges}

Variance-share uncertainty inherits identified-set uncertainty in the mean news
coefficient. Published mean specification B makes the expected and combined
shares set-valued. Diagnostic mean specification A leaves the expected block
constant.

\file{scripts-paper/mean_equation/variance_shares/} forms centered divisor-$T$
second-moment blocks for the expected-SDF and SDF-news PC score blocks and their
covariance.
With $B_E(\theta)$ denoting the expected-SDF coefficient block and $V_c$ the
variance of quarterly real consumption growth, the reported maps are
\begin{align*}
  s_N(\theta) &=
    100\,\frac{\theta\trans S_N\theta}{V_c},\\
  s_E(\theta) &=
    100\,\frac{B_E(\theta)\trans S_EB_E(\theta)}{V_c},\\
  s_{E+N}(\theta) &=
    s_E(\theta)+s_N(\theta)
    +200\,\frac{B_E(\theta)\trans S_{EN}\theta}{V_c}.
\end{align*}
Here $S_N$ and $S_E$ are the divisor-$T$ second-moment matrices for the
SDF-news and expected-SDF PC score blocks. $S_{EN}$ is their cross-moment
matrix, and $V_c$ is the variance of quarterly real consumption growth.

Each reported attained numerical range starts from a feasible lattice with 101
points per axis and expands with feasible values found by multistart local
SLSQP. These additions widen the attained ranges and leave the nonconvex set's
global extrema uncertified. Coherence checks use the contract ratio 0.98
and numerical tolerance $10^{-9}$. Hamilton rounding applies only to the OLS and
tau-zero point columns, making their displayed components sum to displayed block
rows. Set ranges and combined rows remain unreconciled.

\section{Log-variance estimators and set maps}
\label{sec:logvar}

The log-variance analysis maps each admissible mean news coefficient into
residual-scale coefficients on one aligned sample. The log-variance set engine
handles the common numerical search. PPML and the Harvey estimator are the
primary estimators; LAD is dependency-gated.

\subsection{Shared log-variance sample, residual map, and registries}

\src{scripts-paper/log_variance/estimators/log_ols/run.R}{log_var_eq}
constructs the shared log-variance sample. It inner-joins the
published mean-equation rows with \code{l.pc1}--\code{l.pc4} by quarter and
centers the four nominal asset-return PCs without scaling. Define the residual
and log-variance design row by
\[
  e_t(\theta)=W_{1t}-W_{2t}\trans \theta,
  \qquad
  R_t=(1,\mathrm{PC}_{1t},\ldots,\mathrm{PC}_{4t}).
\]
Here $e_t(\theta)$ is the mean-equation residual at mean news coefficient
$\theta$, and $R_t$ contains an intercept and the four centered nominal
asset-return PC scores. Every estimator maps the same mean news coefficient
$\theta\in\Theta(\tau)$ to five residual-scale coefficients $\eta$ with axis
\code{(Intercept), l.pc1, l.pc2, l.pc3, l.pc4}. The implementation under
\file{scripts-paper/log_variance/engine/} is the log-variance set engine. It owns
feasible grids, endpoint scans, endpoint polishing, box-escape checks, cold-start
endpoint comparisons, fit-evaluation budgets, and
stored endpoint-status assembly. An estimator object supplies its fit, Jacobian, domain, and
side-specific hooks. Endpoint polishing uses SLSQP for smooth objectives with a
gradient and derivative-free Constrained Optimization BY Linear Approximations
(COBYLA) otherwise. LAD follows the COBYLA path.

The two-step log-OLS benchmark is the separate foundation. The log-variance
estimator registry's primary order is PPML followed by the Harvey estimator; its
extension entries add LAD when the decision and dependency gate source it. The
log-variance estimator registry is authoritative for dependencies and
capabilities. PPML and the Harvey estimator have \code{set_bootstrap}. LAD does
not. A different object, the bounds-rendering registry
\code{logvar_bounds_tau_registry}, begins with the two-step log-OLS benchmark and
receives PPML, Harvey estimator, and conditional LAD entries in runner order.

\subsection{Two-step log-OLS benchmark}

For every admissible $\theta$, the two-step log-OLS benchmark computes
\[
  \widehat\eta_{\mathrm{logOLS}}(\theta)
  =(R\trans R)^{-1}R\trans\log\!\left(e(\theta)^2\right).
\]
Exact zero residuals are holes in this map's domain; the implementation does
not floor them. The crossing module uses projection weights to determine the
coefficient-specific direction of divergence. A certified crossing may
imply a genuine infinite endpoint. A suspected but uncertified
crossing receives \code{unreliable}, not an artificial finite bound.

The log-OLS analysis stage also creates the shared design, the tau-zero map, and
the first entry in the bounds-rendering registry. It computes refined mean
coefficient boxes immediately afterward because the nonlinear estimators consume
those exact relaxation-value-keyed boxes.

\subsection{Common set map and second search audit design for PPML and the Harvey estimator}
\label{sec:ppml-harvey-search}

PPML quasi-Poisson log-link QMLE and Harvey Gaussian multiplicative-variance
QMLE are the two primary log-variance estimators. At every admissible $\theta$,
both use
the squared-residual response and exponential conditional-variance map
\[
  y_t(\theta)=e_t(\theta)^2,
  \qquad
  \mu_t(\eta)=\exp(R_t\trans\eta).
\]
For each displayed relaxation value $\tau$, PPML and the Harvey estimator share
the log-variance set engine, warm-refined mean coefficient box, search seed, and
four-lag Bartlett HAC standard errors for point columns. These analytic standard
errors supplement the unified bootstrap. Their primary full-sample set maps,
second search audits, and bootstrap set maps use the matched controls in
\src{scripts-paper/config/inference_search_control.R}{PAPER_LOGVAR_BUDGETS} and
\src{scripts-paper/log_variance/estimators/controls.R}{LOGVAR_SEARCH_CONTROL}.
A selected-grid cap limits retained feasible rows. A fit-evaluation budget
limits estimator fits used by the grid scan, endpoint search, and checks. Each
name limits estimator-fit claims within one set-engine search. A separate
control sets the bootstrap replication count.

\begin{table}[htbp]
  \centering
  \caption{Matched Poisson pseudo-maximum-likelihood (PPML) and Harvey search
  settings. Each fit-evaluation budget resets for every relaxation value shown
  or searched. Each start count is the scan-derived pool for one residual-scale
  coefficient and one endpoint side; the feasible search seed and any carried
  endpoint-attaining points enter separately.}
  \label{tab:ppml-harvey-search}
  \small
  \begin{tabular}{@{}>{\raggedright\arraybackslash}p{0.20\linewidth}
                      >{\raggedright\arraybackslash}p{0.36\linewidth}
                      >{\raggedright\arraybackslash}p{0.36\linewidth}@{}}
    \toprule
    Run & PPML & Harvey \\
    \midrule
    Primary full-sample set map &
    At most 4,000 selected feasible-grid rows with deterministic lattice
    coarsening; fresh 20,000-fit-evaluation budget per displayed relaxation value;
    scan-derived start pool of three per coefficient and endpoint side. &
    At most 4,000 selected feasible-grid rows with deterministic lattice
    coarsening; fresh 20,000-fit-evaluation budget per displayed relaxation value;
    scan-derived start pool of three per coefficient and endpoint side. \\
    Second search audit &
    At most 8,000 Morton-selected feasible-grid rows; fresh
    40,000-fit-evaluation budget per displayed relaxation value; scan-derived start pool
    of five per coefficient and endpoint side. &
    At most 4,000 selected feasible-grid rows with deterministic lattice
    coarsening; fresh 40,000-fit-evaluation budget per displayed relaxation value;
    scan-derived start pool of five per coefficient and endpoint side. \\
    Bootstrap set map &
    At most 1,500 selected feasible-grid rows with deterministic lattice
    coarsening; 8,000-fit-evaluation budget per searched relaxation value; scan-derived
    start pool of three per coefficient and endpoint side. &
    At most 1,500 selected feasible-grid rows with deterministic lattice
    coarsening; 8,000-fit-evaluation budget per searched relaxation value; scan-derived
    start pool of three per coefficient and endpoint side. \\
    \bottomrule
  \end{tabular}
\end{table}

The feasible search seed enters every run separately. A primary full-sample set
map may also add endpoint-attaining points from the previous displayed
relaxation value. ``Five starts'' names the same scan-derived pool for
both second search audits. It does not count all optimizer trials.

Table~\ref{tab:ppml-harvey-search} shows identical primary full-sample set map
settings and the larger scan-derived start pool and fit-evaluation budget shared
by the two second search audits. Those audits also share warm-refined mean
coefficient boxes and one fresh set-engine cache, independent of the primary
cache, across the audited relaxation values. Both second search audits enable
the cold-start endpoint comparison; both bootstrap set maps disable it.

After each audit, the same endpoint reconciliation rule keeps the more extreme
certified endpoint and records its source. Endpoint movement beyond the
agreement tolerance, a status mismatch, or a second search audit failure changes
a bounded result to \code{unreliable}.

The second search audits differ in selected-grid and estimator-object handling.
PPML uses 8,000 Morton-selected rows and a new PPML estimator object. The Harvey
estimator retains the 4,000-row deterministic grid and existing estimator
object. The PPML grid-coverage check probes grid coverage; the Harvey second
search audit checks endpoint stability on its original grid. Both remain local
audits of the nonconvex set and provide no global-endpoint certificate.

\subsection{PPML quasi-Poisson log-link QMLE}

\src{scripts-paper/log_variance/estimators/ppml/acceptance.R}{logvar_ppml_accept}
computes and checks the PPML estimating equation
\[
  R\trans\left[e(\theta)^2-\exp(R\eta)\right]=0
\]
and \src{scripts-paper/log_variance/estimators/ppml/fit.R}
{logvar_ppml_fit_response} solves it through \code{glm.fit()} with a
quasi-Poisson family and log link. A pilot freezes the PPML response scale; the
reported intercept restores the log of that scale. Acceptance requires finite
coefficients and fitted means,
convergence without a boundary condition, normalized score at most $10^{-8}$,
and adequate rank and conditioning. \code{LOGVAR_PPML_CONTROL} fixes the fit
start order and pilot behavior. It also fixes the PPML grid-coverage check's
Morton-key precision and exact-double guard. After the primary full-sample set
map, the grid-expanded PPML second search audit applies the shared
second search audit and endpoint reconciliation design in
Section~\ref{sec:ppml-harvey-search}.

\file{log_variance/estimators/ppml/fit.R} also defines the convenience primitive
\code{logvar_ppml_score} and information primitive \code{logvar_ppml_info}.
The production fit does not call either named function; the active acceptance
code evaluates the same score and information cross-products directly. The
primitives are available to tests and future moment code but are not separate
production paths.

The PPML covariance variants are naive, heteroskedasticity-consistent HC0 and
HC1, and HAC. The reporting policy in Table~\ref{tab:contract} selects Bartlett
HAC with four lags.

\subsection{Harvey Gaussian multiplicative-variance QMLE}

\src{scripts-paper/log_variance/estimators/harvey/solver_primitives.R}{hv_eval}
evaluates the Gaussian negative-log-likelihood objective
\[
  Q_H(\eta;\theta)=\frac12\sum_t
    \left[R_t\trans\eta+y_t(\theta)\exp(-R_t\trans\eta)\right].
\]
and the first-order-condition moment $R\trans(y/\mu-1)$, the negative of twice
the score $\frac12R\trans(1-y/\mu)$.
\src{scripts-paper/log_variance/estimators/harvey/likelihood.R}
{logvar_harvey_info} gives observed information
$\frac12R\trans\operatorname{diag}(y/\mu)R$; the corresponding expected
information is $\frac12R\trans R$. The solver combines observed-Newton and
Fisher steps, backtracking, multiple starts, PPML warm starts, and a
linear-program recession certificate. The Harvey estimator uses the same primary
full-sample set map, bootstrap set map, and second search audit structure as
PPML; Table~\ref{tab:ppml-harvey-search} isolates the grid difference.

\file{log_variance/estimators/harvey/likelihood.R} also defines
\code{logvar_harvey_objective}, \code{logvar_harvey_score}, and
\code{logvar_harvey_moment}. The current solver calls \code{hv_eval}, which
implements those formulas directly; it does not call the three named convenience
primitives. They remain available to tests and the source-loaded, driver-blocked
joint-GMM branches.

\src{scripts-paper/log_variance/estimators/harvey/sensitivity_and_reporting.R}
{logvar_harvey_sensitivity_gate} implements the fixed-grid Harvey second search
audit and passes its result to the shared endpoint reconciliation rule. The
TeX term \emph{second search audit} maps to PPML code named
\code{coverage}/\code{coverage_audit} and Harvey code named
\code{gate}/\code{sensitivity_audit}.
In bootstrap tau-zero evaluation, the code reuses the stored
Harvey estimator \code{point_fit}; calling the Harvey estimator's generic
\code{fit_at_b()} would apply a different start policy and could define a
different tau-zero log-variance fit. The Harvey estimator covariance variants
are expected, observed, outer product of gradients (OPG), robust, and HAC. The
reporting policy in Table~\ref{tab:contract} selects four-lag Bartlett HAC.

\subsection{Dependency-gated LAD conditional-median map}

The committed decision in \file{scripts-paper/config/decisions/lad.dcf} is
approved for \code{quantreg} 6.1. The runner first checks whether
the namespace exists, then validates the decision against the installed version.
An approved estimator with an absent or mismatched dependency aborts the run and
prevents a silent skip.

The tracked decision accepts \code{unanswered}, \code{declined}, or
\code{approved}. The gate reader classifies a nonexistent path as
\code{unanswered}, but the runner obtains this record through required
\code{paper_path()} resolution, so deleting the tracked file stops before the
reader. Installation consent
accepts \code{not_asked}, \code{not_needed}, \code{declined}, or
\code{approved}. A coherent edit updates the decision, consent, rationale,
timestamp, and recorded package version together. Only approval with the exact
installed version sources LAD.

When the LAD dependency gate returns \code{source_lad=TRUE}, LAD fits
\[
  2\log|e_t(\theta)|=R_t\trans\eta+v_t.
\]
Here $v_t$ is the median-regression residual. The estimator uses median
regression with \code{quantreg::rq.fit(method="br", tau=0.5)}. The \code{fn} method is a
nonuniqueness probe. Here \code{tau=0.5} is the quantile probability for the
median regression, not the identification relaxation value. Exact-zero
observations remain unchanged and trigger domain failures. The code
distinguishes the attained finite coefficient range from approximate closure
limits and can demote a
near-crossing result to \code{unreliable}. LAD can produce a
bounds-by-relaxation-value figure, table page, fitted conditional-median
absolute-residual envelope, and closure CSV. The unified bootstrap excludes LAD
by design. The decision therefore changes conditional LAD execution and its
three artifact records without invalidating the unified-bootstrap cache.

\subsection{Diagnostics, residual-dynamics gate, and EGARCH-X decision router}
\label{sec:diagnostics}

Diagnostics run at three points. Mean estimation assembles the mean-equation
relevance record. The joint-null witness diagnostic, joint-GMM replication
diagnostic, residual-dynamics gate, and EGARCH-X decision router run before
bootstrap. A publication stage produces the heteroskedasticity diagnostic
battery. Their roles are:
\begin{itemize}
  \item The mean-equation relevance record and heteroskedasticity diagnostic
    battery diagnose identifying variation and leave the mean estimator
    unchanged.
  \item The joint-null witness diagnostic searches for a feasible $\theta$ whose
    four nominal asset-return-PC slopes are jointly zero. It reports a compatible
    joint-null witness or that compatibility was not demonstrated. It is not a
    hypothesis test, so rejection and p-value language would be incorrect.
  \item The current joint-GMM replication diagnostic is confined to graph,
    algebra, and software invariance. It supplies no extra identifying
    restriction or overidentification test. Although its tracked record contains
    future ratification fields, the current driver accepts only the no-answer
    default and stops on every nondefault choice. The record is consumed, but
    those fields cannot enable their substantive branches.
  \item The residual-dynamics gate applies a lag-4 Ljung--Box rule at
    $\alpha_{\mathrm{dyn}}=0.05$ to tau-zero residuals. It supplies the verdict read by the
    EGARCH-X decision router. Estimator choice remains fixed elsewhere.
\end{itemize}

The EGARCH-X decision router validates the residual-dynamics gate record and
rewrites the terminal EGARCH status record. The current runner wires no EGARCH-X
artifact producer and calls the final conditional-route reconciler with
\code{egarch_producer_ran=FALSE}. Final reconciliation therefore requires all
four conditional EGARCH artifact paths to be absent. The
committed decision record binds a
lag-4, $\alpha_{\mathrm{dyn}}=0.05$ \code{non_reject} gate and records both downstream decisions
as \code{not_asked}; a residual-dynamics gate record that does not match that
decision record cannot silently reuse it.

The two approval fields become answerable only after a rejecting fresh gate.
Even then, the current runner sources no dynamic estimator module and reports
\code{egarch_producer_ran=FALSE}. The present graph therefore cannot turn an
approval into EGARCH-X estimates or artifacts; the fields describe a future
route, not a current user control.

\section{Unified bootstrap and inference}
\label{sec:bootstrap}

\subsection{Unified bootstrap stage and two scientific result branches}

\src{scripts-paper/inference/run_bootstrap_stage.R}{run_bootstrap_stage} is the
only active stage that constructs both the primary and volatility-sensitivity
resample-index families. It computes or reuses one payload containing the
published mean, log-variance, and volatility-sensitivity results and manages the
unified bootstrap cache artifact. The paired B-versus-A mean-specification
comparison reuses the primary resample-index family without creating another
resampling protocol or unified bootstrap cache artifact.

The shared bootstrap frame retains every row of the gapless mean sample and
joins nominal asset-return PCs by quarter. A primary bootstrap replication has
the following exact nesting:
\begin{enumerate}
  \item index the shared bootstrap frame with one prebuilt circular-MBB index draw;
  \item estimate published mean specification B once;
  \item construct tau-zero and displayed-relaxation-value geometry once;
  \item evaluate the mean tau-zero point and positive-relaxation-value endpoints
    from that shared state;
  \item apply log-variance complete cases, recenter the resampled nominal
    asset-return PCs,
    and leave their scales unchanged;
  \item evaluate PPML, then the Harvey estimator, using the log-variance
    estimator registry's dependency graph;
  \item record direct tau-zero fits and positive-relaxation-value endpoint records.
\end{enumerate}
Shared-estimation failure affects both branches. Once shared state exists,
result assembly records mean and estimator-specific failures within their own
branches. The full-data volatility anchor uses the same evaluation callback as
a draw.

The primary resample-index family uses $B=10{,}000$, seed 20260708, and circular moving blocks
of length
\[
  \ell=\left\lceil1.5T^{1/3}\right\rceil.
\]
Here $T$ is the number of rows in the shared bootstrap frame.
The volatility-sensitivity resample-index family uses the same replication count
and block length $2\ell$. RNG kinds are Mersenne-Twister, Inversion, and
Rejection. Before parallel work, the builder constructs both resample-index
families, preserves their draw order, records index and post-index RNG hashes,
and restores ambient RNG state. These steps make serial and parallel resampling
identities independent of solver-side randomness.

Transport failures may not exceed $\lfloor0.25B\rfloor$ draws. Separately, for
each log-variance estimator and relaxation value, the corresponding cell must
keep the pooled \code{failed}-status share
across result draws, coefficients, and sides at or below 0.25. A
volatility-sensitivity result draw re-estimates the log-variance branch under the
volatility-sensitivity resample-index family; it is not a rescaling of a primary
standard error.

\subsection{Stored four-status endpoint vocabulary}

\src{scripts-paper/support/identification/status_contract.R}
{PAPER_ENDPOINT_STATUS} defines a closed four-value vocabulary. The ordering
matters because worst-status aggregation follows the entries from top to
bottom.
\begin{table}[htbp]
\centering
\caption{Stored endpoint-status vocabulary. Ordinary solvers emit
\code{bounded}, \code{unbounded}, or \code{unreliable}; result collectors may
add \code{failed}. The order is also the worst-status aggregation order.
\label{tab:endpoint-status}}
\begin{tabular}{@{}lp{0.68\textwidth}@{}}
\toprule
\code{bounded} & A valid finite endpoint has its context-specific numerical certificate.\\
\code{unbounded} & A numerical certificate establishes an open direction.\\
\code{unreliable} & An attempted evaluation lacks a valid certificate.\\
\code{failed} & Collection or a wholesale estimator branch failed.\\
\bottomrule
\end{tabular}
\end{table}
Table~\ref{tab:endpoint-status} implies that a finite number is publishable only
when its status is \code{bounded}.

At $\tau=0$, \code{point} and \code{point_status} are authoritative. The mean
tau-zero point comes from the linear system. The unified bootstrap directly
evaluates each log-variance estimator at that mean tau-zero point. Tau-zero \code{unbounded} is
invalid. Compatibility fields named lower and upper exactly mirror the point and
status for failure accounting; they are not two tau-zero endpoint estimates. At
positive relaxation values, the two endpoint sides and their statuses are
independent.

\subsection{\TauP{} inference at positive relaxation values}

For each full-sample side certified \code{bounded}, the scale is the MAD of
finite bounded bootstrap endpoints on that side. A side needs at least
$\lceil0.50B\rceil$ such draws, a bounded share of at least 0.85 among nonfailed
draws, and a positive finite scale. Failed draws leave the denominator;
unbounded and unreliable draws remain in it but are not usable.

For two \code{bounded} full-sample sides, calibration requires at least
$\lceil0.50B\rceil$ replications with both endpoints bounded and pairs those
endpoints within replication. If one side is \code{bounded} and the other is
\code{unbounded}, it calibrates the single bounded side. Two unbounded sides
suppress the interval; a full-sample \code{unreliable} or \code{failed} status
also suppresses it. A requested relaxation value that fails the endpoint
regularity gate---minimum count, status stability, or positive scale---stops
publication.

The empirical critical value is the
$\lceil(n+1)(1-\alpha_{\mathrm{inf}})\rceil$ order statistic, capped at the
sample maximum. Here $n$ is the number of usable bootstrap statistics and
$\alpha_{\mathrm{inf}}=0.10$ is the inference level.
\TauP{} parameterizes the truth's location by $\lambda\in[0,1]$ and solves the
resulting worst-location problem with a certified Lipschitz branch-and-bound to
tolerance $10^{-4}$. When the cell's \code{reason} is \code{reported},
\code{ci_lower} and \code{ci_upper} are the published \TauP{} confidence
limits; otherwise those fields are unavailable. The containment target \TauS{}, normal approximations,
Imbens--Manski, and Stoye calculations remain diagnostics.

\subsection{Tau-zero bootstrap point/MAD statistic and stars}

When the tau-zero reporting gate passes, the reported tau-zero bootstrap
point/MAD statistic divides the full-sample point by the MAD of its finite
bounded point draws. Its bootstrap absolute-deviation p-value compares
$|\widehat\phi_b^*-\widehat\phi|$ with $|\widehat\phi|$ and uses a plus-one
correction. Here $\widehat\phi$ is the full-sample tau-zero point and
$\widehat\phi_b^*$ is replication $b$'s bounded point draw. Dividing both sides
by the same MAD leaves the comparison unchanged, so the bootstrap
absolute-deviation p-value is unstudentized. The normal p-value may appear only
in diagnostics. A failed gate leaves the statistic, bootstrap p-value, and stars
unavailable and records the reason. Eligible tau-zero table cells report the
bootstrap point/MAD statistic and stars rather than a bootstrap point interval.
Table-facing code calls the statistic a bootstrap t statistic; that label does
not change the p-value construction.

\subsection{Paired B-versus-A mean-specification comparison}

After the unified bootstrap stage,
\src{scripts-paper/mean_equation/inference/spec_comparison.R}
{mean_spec_comparison_draws} re-estimates diagnostic mean specification A on the
exact primary resample-index family already used for published mean
specification B. The analysis stage checks the resample-index family hash,
assembles both specifications, and writes the diagnostic CSV. The paired
B-versus-A mean-specification comparison does not create a second resampling
protocol or a separate unified bootstrap cache artifact.

\section{Unified bootstrap cache identity and reproducibility}
\label{sec:cache}

The unified bootstrap cache reuses expensive draws only when their scientific
identity is unchanged. Reuse is all-or-nothing. Semantic provenance invalidates
stale draws. A change confined to the presentation manifest leaves draws valid;
a presentation-only field stored in a draw-hashed file or whole control object
still invalidates reuse under the current implementation.

The unified bootstrap cache artifact is \code{bootstrap_stage_draws.rds}. Its
bootstrap cache schema 4 payload contains exactly \code{anchor}, \code{mean},
\code{volatility_primary},
\code{volatility_primary_n_failed}, \code{volatility_sensitivity},
\code{volatility_sensitivity_n_failed}, and \code{provenance}.

Semantic provenance compares the resampler, sample size, replication and block
designs, seed, RNG kinds, axes, cache schema, input hash, draw-spec hash,
draw-code hash, runtime hash, both resample-index family hashes, and both
post-index RNG hashes. The runtime hash adds R and platform details, versions of
\code{hetid}, \code{nloptr}, and \code{sandwich}, repository R sources, loaded
namespace functions, Basic Linear Algebra Subprograms (BLAS), Linear Algebra
Package (LAPACK), and scientific controls. The provenance record stores
\code{presentation_sha} for audit only. A presentation helper that changes a
draw estimate must enter the bootstrap draw-code manifest.

Cache invalidation follows source identity at file and whole-object granularity.
The draw-code manifest includes the analysis contract, inference-search control,
log-variance estimator registry and accessors, reporting policy, typed-artifact
code, the full identification directory, the PPML and Harvey directories, and
the unified-bootstrap implementation. The runtime hash also includes the whole
analysis-contract, inference-search, reporting, and serialization control
objects.
Therefore, an edit in one of those authorities invalidates reuse even when the
edited leaf has only a reporting effect. Figure-style and figure-rendering
controls are outside the draw identity. The presentation manifest is recorded
but not compared for reuse.

In \code{reuse} mode, a missing or semantically stale unified bootstrap cache
artifact warns and triggers a complete rerun. In \code{rerun} mode, computation
is unconditional. The promotion routine writes to a temporary location,
performs exact object RDS readback and cache reuse validation, and then replaces
the current artifact. The stage removes legacy separate mean and log-variance
cache names only after the unified bootstrap cache artifact passes validation.
Exact object RDS readback requires object identity at write time; byte identity
across R versions or machines lies outside the contract.

\section{Publication stages}
\label{sec:publication}

\subsection{Inference tables and estimator pages}

The combined-table artifact producer publishes the mean-equation and PPML
inference panels as one combined mean-equation and PPML inference table triple.
The estimator-pages artifact producer publishes a log-variance estimator-pages
table triple assembled from the available result objects. It always emits PPML
and two-step log-OLS benchmark pages, then adds a Harvey estimator page when its
result exists and a LAD page when its dependency-gated result exists. Under the
current runner, the Harvey estimator result exists before this artifact producer
runs. \file{scripts-paper/config/artifact_latex.R} controls table labels and
LaTeX publication registry membership.

All tau-zero stars in these tables come from the bootstrap point/MAD statistic.
Intervals at positive relaxation values use \TauP{}. The artifact producers
treat stored endpoint status as data. Only \code{bounded} status permits a
finite displayed endpoint, regardless of populated numeric slots.

\subsection{Coefficient bounds and fitted-scale figures}

\src{scripts-paper/log_variance/figures/render_bounds_by_tau.R}
{logvar_bounds_tau_entry} writes one figure for every entry in
the bounds-rendering registry \code{logvar_bounds_tau_registry}. The entries
include the separate two-step log-OLS benchmark and each available log-variance
estimator with a set map. The baseline fitted-volatility analysis stage writes fitted
conditional-residual-volatility envelopes for PPML and the Harvey estimator at
the contract baseline. The conditional LAD analysis stage separately writes a
fitted conditional-median absolute-residual envelope when its dependency-gated
result exists. The fitted-volatility sweep over relaxation values first
filters $\{0.05,0.10,0.20,0.30,0.40\}$ against the estimated boundedness
frontier and stops if no value survives. Otherwise it writes one figure per
surviving relaxation value and estimator.
If more than one relaxation value survives, it also writes linear, log-scale,
lower-normalized, and upper-normalized combined panels.

The manifest marks all 18 PPML and Harvey estimator artifact records for the
fitted-volatility sweep over relaxation values \code{required}; runtime frontier
filtering may omit paths. Artifact completeness and render acceptance therefore
require an independent missing-file check after any run that reaches the end.

The mean-region analysis stages prepare shared refined geometry and render two
projection figures plus 16 three-dimensional regions spanning four relaxation
values, two unit systems, and two OLS-projection modes.

\subsection{Heteroskedasticity diagnostic battery}

The late heteroskedasticity analysis stage recomputes the heteroskedasticity
diagnostic battery and publishes its three-panel table triple. The battery
includes White, Breusch--Pagan,
Goldfeld--Quandt, the Harvey heteroskedasticity test, regime-dependent Anscombe,
Glejser, the Breusch--Pagan Lagrange-multiplier (LM) statistic,
first-order autoregressive-conditional-heteroskedasticity [ARCH(1)], and joint
relevance diagnostics. The battery diagnoses identifying variation; estimator
selection remains fixed elsewhere.

\subsection{SDF approximation-error variance bounds}

At each of the 39 quarterly maturity nodes, $i=3,6,\ldots,117$, the pipeline
computes one SDF approximation-error variance bound for each reported series.
These are plug-in estimates of population bounds rather than finite-sample
guarantees. The package envelope arm for SDF news is
\[
  U_i=\frac14\widehat c_i(\widehat k_i+\widehat k_{2i}).
\]
Here $\widehat c_i$ is the sample envelope. $\widehat k_i$ is the
realized-forecast-error fourth moment. $\widehat k_{2i}$ is the price-news
fourth moment. \src{R/compute_news_q_bound.R}{compute_news_q_bound} supplies the
q-arm. The reported SDF-news approximation-error variance-bound series is the
pointwise minimum of these two arms. The q-arm retains the full
first-order-cancelled remainder. The envelope arm keeps only the leading
fourth-order term.

\src{R/compute_expected_sdf_variance_bound.R}{compute_expected_sdf_variance_bound}
produces the pointwise-minimum expected-SDF approximation-error variance-bound
series.
Its q-arm is exact in the same sense. Its component arm is a leading fourth-order
term. Whenever an envelope or component arm is the minimum,
the reported value inherits that leading-term qualification. All arms use the
validated ACM yields and term premia. Reported series must be finite and
strictly positive; the news q-arm alone may use documented positive infinity as
a conservative fallback before the pointwise minimum.

Before either variance-bound artifact producer, the quoted-number reconstruction
and assertion module reconstructs the claims used by the paper and fails if they
drift outside source-coded tolerances. The
current assertions include a maximum news bound near $4.19\times10^{-9}$ and a
maximum expected-SDF bound near $3.19\times10^{-9}$. They also require the
implied standard deviations to remain below 0.7 basis points and below 3 percent
of a typical quarterly yield move. The asserted typical one-quarter-yield move
is near 0.00239. The module then writes a typed CSV and a Markdown report. After
all assertions pass, the runner writes the SVG and the summary-table triple.

\subsection{Descriptive-statistics report}

The last substantive artifact producer full-joins quarterly real consumption
growth, the three SDF-PC score blocks, selected expected-SDF and SDF-news
maturities, and selected realized quarterly yield volatilities.

The artifact producer reports missing date--variable combinations; writes
separate summary, correlation, and OLS benchmark tables; creates a two-page PDF
of time series and histograms; and compiles the wrapper TeX into the report PDF.
Correlations use pairwise complete observations, and summary rows omit missing
values.

\subsection{Final conditional-route reconciliation and cleanup}

The final conditional-route reconciler separates permission from producer
execution. A conditional artifact producer may report \code{ran=TRUE} only if it
has \code{requested=TRUE} and all artifact paths with its lifecycle status
exist. A route
with \code{ran=FALSE} must leave every associated artifact path absent.
Requested but not run is a valid recorded outcome. The resulting schema 1.1.0
conditional-route state record contains routes, audits of the startup cleanup
of conditional artifacts, and per-path presence. The final conditional-route
reconciler writes the record
through exact object RDS readback.

Reconciliation covers the seven conditional artifact records; a separate
acceptance check must establish existence and freshness for the 77 required
records. Final LaTeX-sidecar cleanup is the runner's last statement, so a cleanup
failure may occur after every scientific artifact producer and the
conditional-route state record write.

\section{Artifact manifest}
\label{sec:artifacts}

The artifact manifest defines the expected artifact universe, paths, ownership,
and lifecycle status. Required-file existence and freshness require runtime
acceptance.

\subsection{Manifest shape and path rule}

The expanded artifact manifest contains 84 unique artifact records: 50 literal
records, 18 fitted-volatility-sweep records, and 16 generated mean-region
records. Of these, 77 are \code{required}, three \code{conditional_lad}, and four
\code{conditional_egarch}; together they occupy seven output groups. Every
current artifact path, stored in \code{new_path}, follows
\begin{center}
  \file{scripts-paper/output/}\code{<group>/<basename>}.
\end{center}
The retained \code{old_path} stores the legacy flat
\file{scripts-paper/output/}\code{<basename>} location.

\begin{table}[htbp]
\centering
\small
\caption{Counts for all 84 expanded artifact-manifest records by output group.
LAD means least absolute deviations. EGARCH-X means exponential generalized
autoregressive conditional heteroskedasticity with exogenous regressors. The
manifest contains 77 required records, three records conditional on LAD, and
four records conditional on EGARCH-X. Counts describe expected records, not
current file presence.
\label{tab:artifact-groups}}
\begin{tabular}{@{}lr@{}}
\toprule
Group & Records\\
\midrule
\file{tables/descriptive_statistics} & 3\\
\file{tables} & 15\\
\file{figures} & 46\\
\file{figures/descriptive_statistics} & 1\\
\file{reports} & 3\\
\file{diagnostics} & 11\\
\file{state} & 5\\
\midrule
Total & 84\\
\bottomrule
\end{tabular}
\end{table}

Table~\ref{tab:artifact-groups} shows that figures are the largest output group,
with 46 of the 84 artifact records.

Artifact lifecycle status is not execution reachability. The exact crosswalk is:
\begin{itemize}
  \item \code{required} means the manifest expects the record regardless of a
    conditional route. Its producer stage is always attempted if earlier stages
    succeed. The fitted-volatility sweep stops if no configured value survives;
    with one or more surviving values, runtime filtering can still omit required
    sweep paths. No final gate checks all 77 required paths.
  \item \code{conditional_lad} means the file must exist exactly when the LAD
    producer ran. LAD is source-selected subject to exact
    \code{quantreg} 6.1; a declined or unanswered record selects absence.
  \item \code{conditional_egarch} means a future EGARCH-X producer would own the
    file. The current runner has no such producer, deletes these paths at startup,
    and requires them to remain absent.
\end{itemize}
Thus the R, L, and E codes below classify artifact obligations. They do not mean
always attempted, selected path, or available but not called.

\begin{landscape}
\subsection{The 50 literal records}

The 50 literal artifact records assign one ID to each manifested file. A
producer may write a publication triple or paired diagnostic formats; the
distinct artifact IDs keep each file separately addressable. The artifact writer
determines typed behavior. The artifact ID only addresses the file.

\begingroup
\footnotesize
\setlength{\tabcolsep}{3pt}
\begin{longtable}{@{}>{\raggedright\arraybackslash}p{0.28\textwidth}
  >{\raggedright\arraybackslash}p{0.33\textwidth}
  >{\raggedright\arraybackslash}p{0.32\textwidth}c@{}}
\caption{The 50 literal artifact records. Group/basename paths are relative to
\file{scripts-paper/output/}; producer paths are relative to
\file{scripts-paper/}. Status R means required, L means conditional least
absolute deviations (LAD), and E means conditional exponential generalized
autoregressive conditional heteroskedasticity with exogenous regressors
(EGARCH-X).\label{tab:literal-artifacts}}\\
\toprule
ID & Group/basename & Producer under \file{scripts-paper/} & Status\\
\midrule
\endfirsthead
\toprule
ID & Group/basename & Producer under \file{scripts-paper/} & Status\\
\midrule
\endhead
\code{summary_statistics_table} & \file{tables/descriptive_statistics/summary_stats.tex} & \file{reports/build_descriptive_statistics.R} & R\\
\code{correlations_table} & \file{tables/descriptive_statistics/correlations.tex} & same descriptive producer & R\\
\code{ols_mean_equation_table} & \file{tables/descriptive_statistics/ols_mean_eq.tex} & same descriptive producer & R\\
\code{heteroskedasticity_table} & \file{tables/hetero_tests.tex} & \file{mean_equation/diagnostics/heteroskedasticity/render_table.R} & R\\
\code{heteroskedasticity_standalone_tex} & \file{tables/hetero_tests_standalone.tex} & same heteroskedasticity producer & R\\
\code{heteroskedasticity_standalone_pdf} & \file{tables/hetero_tests_standalone.pdf} & same heteroskedasticity producer & R\\
\code{structural_var_inference_table} & \file{tables/structural_var_inference.tex} & \file{log_variance/tables/render_combined_inference_table.R} & R\\
\code{structural_var_inference_standalone_tex} & \file{tables/structural_var_inference_standalone.tex} & same combined-table producer & R\\
\code{structural_var_inference_standalone_pdf} & \file{tables/structural_var_inference_standalone.pdf} & same combined-table producer & R\\
\code{structural_var_estimators_table} & \file{tables/structural_var_estimators.tex} & \file{log_variance/tables/render_estimator_pages.R} & R\\
\code{structural_var_estimators_standalone_tex} & \file{tables/structural_var_estimators_standalone.tex} & same estimator-page producer & R\\
\code{structural_var_estimators_standalone_pdf} & \file{tables/structural_var_estimators_standalone.pdf} & same estimator-page producer & R\\
\code{variance_share_table} & \file{tables/var_share.tex} & \file{mean_equation/variance_shares/render_variance_share_table.R} & R\\
\code{variance_share_standalone_tex} & \file{tables/var_share_standalone.tex} & same variance-share producer & R\\
\code{variance_share_standalone_pdf} & \file{tables/var_share_standalone.pdf} & same variance-share producer & R\\
\code{mean_bounds_figure} & \file{figures/set_id_bounds_tau.svg} & \file{mean_equation/inference/compute_bounds_by_tau.R} & R\\
\code{mean_projections_figure} & \file{figures/set_id_projections_sd.svg} & \file{mean_equation/figures/render_projections.R} & R\\
\code{mean_projections_ols_figure} & \file{figures/set_id_projections_sd_ols.svg} & same projection producer & R\\
\code{log_ols_bounds_figure} & \file{figures/log_var_eq_bounds_tau_logols.svg} & \file{log_variance/figures/render_bounds_by_tau.R} & R\\
\code{ppml_bounds_figure} & \file{figures/log_var_eq_bounds_tau_ppml.svg} & same bounds-figure producer & R\\
\code{harvey_bounds_figure} & \file{figures/log_var_eq_bounds_tau_harvey.svg} & same bounds-figure producer & R\\
\code{lad_bounds_figure} & \file{figures/log_var_eq_bounds_tau_lad.svg} & same bounds-figure producer & L\\
\code{ppml_fitted_volatility_figure} & \file{figures/log_var_eq_fitted_volatility_ppml.svg} & \file{log_variance/figures/fitted_volatility/run.R} & R\\
\code{harvey_fitted_volatility_figure} & \file{figures/log_var_eq_fitted_volatility_harvey.svg} & same fitted-volatility producer & R\\
\code{lad_fitted_volatility_figure} & \file{figures/log_var_eq_fitted_volatility_lad.svg} & \file{log_variance/figures/fitted_volatility/run_lad.R} & L\\
\code{descriptive_figures} & \file{figures/descriptive_statistics/figures.pdf} & \file{reports/build_descriptive_statistics.R} & R\\
\code{descriptive_report_tex} & \file{reports/descriptive_stats.tex} & same descriptive producer & R\\
\code{descriptive_report_pdf} & \file{reports/descriptive_stats.pdf} & same descriptive producer & R\\
\code{mean_inference_diagnostics} & \file{diagnostics/set_id_inference_diagnostics.csv} & \file{inference/run_bootstrap_stage.R} & R\\
\code{mean_spec_comparison} & \file{diagnostics/set_id_spec_comparison.csv} & \file{mean_equation/inference/spec_comparison.R} & R\\
\code{log_variance_inference_diagnostics} & \file{diagnostics/log_var_eq_set_inference_diagnostics.csv} & \file{inference/run_bootstrap_stage.R} & R\\
\code{joint_null_csv} & \file{diagnostics/log_var_eq_joint_null.csv} & \file{log_variance/diagnostics/joint_null/run.R} & R\\
\code{joint_null_rds} & \file{diagnostics/log_var_eq_joint_null.rds} & same joint-null producer & R\\
\code{joint_gmm_csv} & \file{diagnostics/log_var_eq_joint_gmm.csv} & \file{log_variance/diagnostics/joint_gmm/pipeline_driver.R} & R\\
\code{joint_gmm_rds} & \file{diagnostics/log_var_eq_joint_gmm.rds} & same joint-GMM producer & R\\
\code{lad_closure_diagnostics} & \file{diagnostics/log_var_eq_lad_closure.csv} & \file{log_variance/estimators/lad/run_sets.R} & L\\
\code{bootstrap_stage_draws} & \file{state/bootstrap_stage_draws.rds} & \file{inference/run_bootstrap_stage.R} & R\\
\code{dynamics_gate} & \file{state/log_var_eq_dynamics_gate.rds} & \file{log_variance/diagnostics/dynamics/run_gate.R} & R\\
\code{egarch_status} & \file{state/log_var_eq_egarch_status.rds} & \file{log_variance/diagnostics/dynamics/run_gate.R} and \file{log_variance/extensions/egarch/run_route.R} & R\\
\code{conditional_route_status} & \file{state/conditional_route_status.rds} & \file{run_pipeline.R} & R\\
\code{egarch_pilot_state} & \file{state/log_var_eq_egarch_pilot.rds} & \file{log_variance/extensions/egarch} & E\\
\code{egarch_results_csv} & \file{diagnostics/log_var_eq_egarch_x.csv} & same EGARCH extension & E\\
\code{egarch_results_rds} & \file{diagnostics/log_var_eq_egarch_x.rds} & same EGARCH extension & E\\
\code{egarch_bounds_figure} & \file{figures/log_var_eq_bounds_tau_egarch_x.pdf} & same EGARCH extension & E\\
\code{variance_bound_figure} & \file{figures/variance_bounds.svg} & \file{variance_bounds/figures/render_bounds.R} & R\\
\code{variance_bound_summary_table} & \file{tables/variance_bounds_summary.tex} & \file{variance_bounds/tables/render_summary_table.R} & R\\
\code{variance_bound_summary_standalone_tex} & \file{tables/variance_bounds_summary_standalone.tex} & same variance-bound table producer & R\\
\code{variance_bound_summary_standalone_pdf} & \file{tables/variance_bounds_summary_standalone.pdf} & same variance-bound table producer & R\\
\code{quoted_numbers_csv} & \file{diagnostics/approximation_error_quoted_numbers.csv} & \file{variance_bounds/quoted/run.R} & R\\
\code{quoted_numbers_md} & \file{reports/approximation_error_quoted_numbers.md} & same quoted-number producer & R\\
\bottomrule
\end{longtable}
\endgroup
\end{landscape}

Table~\ref{tab:literal-artifacts} shows that lifecycle attaches to records,
allowing required and conditional files to share an output group.

\subsection{Artifact records for the fitted-volatility sweep over relaxation values}

The manifest builder generates 18 records for the fitted-volatility sweep over
relaxation values from two estimators and nine exact tails. For estimator $e$
and tail $t$, the ID is \code{<e>_fitted_vol_<t>}. The variant is
\code{<e>_<t>}. The family is \code{fitted_volatility_sweep}. All 18 artifact
records are required.
\file{log_variance/figures/fitted_volatility/run_tau_sweep.R} is the registered
producer for all 18 records, but its current body does not guarantee all 18. It
writes scalar records only for configured relaxation values below the runtime
boundedness frontier and aggregate records only when more than one value
survives. It stops if no configured value survives. Scalar tails use
\code{tau0p05}, \code{tau0p1}, \code{tau0p2}, \code{tau0p3}, and
\code{tau0p4}. The four aggregate tails are \code{combined},
\code{combined_log}, \code{combined_lower_normalized}, and
\code{combined_upper_normalized}.

The basename formula is
\begin{center}
  \code{figures/log_var_eq_fitted_volatility_<e>_<f(t)>.svg},
\end{center}
Here $f(t)=t$ for a scalar tail. For an aggregate tail, $f(t)$ replaces the
leading \code{combined} with \code{tau_combined}. The deterministic prefix is
common to every record.

\begin{table}[htbp]
\centering
\small
\caption{The 18 required artifact records for the fitted-volatility sweep over
relaxation values. PPML means Poisson pseudo-maximum likelihood.
Each estimator has the five scalar and four aggregate tails shown. IDs use
\code{<estimator>_fitted_vol_<ID tail>}; paths use
\file{figures/log_var_eq_fitted_volatility_<estimator>_<basename tail>.svg}.
\label{tab:sweep-artifacts}}
\begin{tabular}{@{}>{\raggedright\arraybackslash}p{0.18\textwidth}
  >{\raggedright\arraybackslash}p{0.36\textwidth}
  >{\raggedright\arraybackslash}p{0.36\textwidth}@{}}
\toprule
Estimator & ID tail after \code{<estimator>_fitted_vol_} & Basename tail\\
\midrule
PPML & \code{tau0p05}, \code{tau0p1}, \code{tau0p2}, \code{tau0p3}, \code{tau0p4} & same five tails\\
PPML & \code{combined} & \code{tau_combined}\\
PPML & \code{combined_log} & \code{tau_combined_log}\\
PPML & \code{combined_lower_normalized} & \code{tau_combined_lower_normalized}\\
PPML & \code{combined_upper_normalized} & \code{tau_combined_upper_normalized}\\
Harvey estimator & \code{tau0p05}, \code{tau0p1}, \code{tau0p2}, \code{tau0p3}, \code{tau0p4} & same five tails\\
Harvey estimator & \code{combined} & \code{tau_combined}\\
Harvey estimator & \code{combined_log} & \code{tau_combined_log}\\
Harvey estimator & \code{combined_lower_normalized} & \code{tau_combined_lower_normalized}\\
Harvey estimator & \code{combined_upper_normalized} & \code{tau_combined_upper_normalized}\\
\bottomrule
\end{tabular}
\end{table}

Table~\ref{tab:sweep-artifacts} shows exact manifest symmetry: PPML and the
Harvey estimator each receive five scalar and four aggregate artifact records.

For example, the first PPML ID is \code{ppml_fitted_vol_tau0p05}. Its path is
\file{figures/log_var_eq_fitted_volatility_ppml_tau0p05.svg}. The first
aggregate Harvey ID is \code{harvey_fitted_vol_combined}. Its path is
\file{figures/log_var_eq_fitted_volatility_harvey_tau_combined.svg}. These
formulas exhaust the Cartesian product used by the artifact manifest builder.

\subsection{Generated mean-region artifact records}

The manifest builder generates 16 required region records from four relaxation
values, two unit systems, and two OLS modes. They have no artifact lookup
family/variant; artifact producers resolve them by artifact ID.
\file{mean_equation/figures/render_region_3d.R} produces all 16. The paper
consumes them, and each has required artifact lifecycle status.
The exact artifact ID and basename formulas cover all combinations.

\begin{table}[htbp]
\centering
\small
\caption{The 16 required three-dimensional mean-region records. In each row,
\code{<t>} takes \code{tau0p05}, \code{tau0p1}, \code{tau0p2}, and
\code{tau0p3}; every basename lies under \file{figures/}. ``SD'' means
standard-deviation units, and OLS means ordinary least squares. Each row
represents four records.
\label{tab:region-artifacts}}
\begin{tabular}{@{}p{0.13\textwidth}p{0.14\textwidth}p{0.66\textwidth}@{}}
\toprule
OLS mode & Units & Four IDs and basenames\\
\midrule
None & SD & \code{mean_region_sd_<t>} / \code{set_id_region_3d_sd_<t>.svg}\\
None & coefficient & \code{mean_region_b_<t>} / \code{set_id_region_3d_b_<t>.svg}\\
Projected & SD & \code{mean_region_ols_projected_sd_<t>} / \code{set_id_region_3d_ols_projected_sd_<t>.svg}\\
Projected & coefficient & \code{mean_region_ols_projected_b_<t>} / \code{set_id_region_3d_ols_projected_b_<t>.svg}\\
\bottomrule
\end{tabular}
\end{table}

Table~\ref{tab:region-artifacts} shows that OLS mode changes the artifact ID
without changing the relaxation-value or unit-system axes.
\subsection{Artifact lookup families and publication triples}

Only three artifact lookup families exist:
\begin{itemize}
  \item \code{logvar_bounds_tau}: \code{logols}, \code{ppml},
    \code{harvey}, conditional \code{lad}, and conditional \code{egarch_x};
  \item \code{fitted_volatility}: \code{ppml}, \code{harvey}, and conditional
    \code{lad};
  \item \code{fitted_volatility_sweep}: the 18 estimator/tail variants in
    Table~\ref{tab:sweep-artifacts}.
\end{itemize}
These are artifact lookup families, not estimand classes. The LAD variant under
\code{fitted_volatility} is a fitted conditional-median absolute-residual
envelope, whereas the variants for PPML and the Harvey estimator are fitted
conditional-residual-volatility envelopes.
An artifact lookup family/variant query must resolve exactly one artifact record.
The other 58 artifact records have blank family and variant fields and are
addressed by artifact ID.

The LaTeX publication registry contains five triples. Each triple binds a
fragment, standalone TeX source, and standalone PDF.
\begin{table}[htbp]
\centering
\scriptsize
\caption{LaTeX publication-registry triples. Each code literal is an artifact
ID in the manifest. Registry assertions check membership and semantic labels;
they do not compile or inspect the files. PDF means Portable Document Format.
\label{tab:latex-publications}}
\begin{tabular}{@{}p{0.28\textwidth}p{0.32\textwidth}p{0.32\textwidth}@{}}
\toprule
Fragment & Standalone TeX & Standalone PDF\\
\midrule
\code{heteroskedasticity_table} & \code{heteroskedasticity_standalone_tex} & \code{heteroskedasticity_standalone_pdf}\\
\code{structural_var_estimators_table} & \code{structural_var_estimators_standalone_tex} & \code{structural_var_estimators_standalone_pdf}\\
\code{structural_var_inference_table} & \code{structural_var_inference_standalone_tex} & \code{structural_var_inference_standalone_pdf}\\
\code{variance_share_table} & \code{variance_share_standalone_tex} & \code{variance_share_standalone_pdf}\\
\code{variance_bound_summary_table} & \code{variance_bound_summary_standalone_tex} & \code{variance_bound_summary_standalone_pdf}\\
\bottomrule
\end{tabular}
\end{table}
Table~\ref{tab:latex-publications} makes registry membership file-specific: each
publication occupies three independently addressable artifact IDs.

\section{Validation boundaries}
\label{sec:validation}

The production runner invokes none of the validation or test entrypoints in this
section. \file{scripts-paper/tests/run_tests.R}, the two-root published-table
comparator, and \file{docs/quality-check.R} are separate entrypoints. Typed writer
readback checks are different: a production writer performs its own readback
when that writer is reached.

\subsection{What source-level checks establish}

Artifact-manifest assembly rejects unused producer or consumer codes and invalid
artifact-group or lifecycle-status codes. Public artifact configuration rejects
missing values, duplicate artifact IDs or basenames, incomplete artifact lookup
family/variant pairs, and duplicate family/variant coordinates. The topology
checks parse and style-check the R tree, inspect source references and
source-once behavior, and validate path topology. Artifact completeness lies
outside their scope.

Typed CSV writers check a round trip at the configured numerical tolerance.
Exact-RDS writers require \code{identical()} after readback. The published-table
projection comparison checks relative table paths, numerical coordinates, token
counts, values under a precision-aware rounding rule, and immediately attached
stars. Its scope excludes captions, labels, notes, other nonnumeric cells, LaTeX
structure, rendered layout, and non-table artifacts. The topology, round-trip,
readback, and projection checks establish only their named properties.

\subsection{Limits of source-level checks}

The current topology checks use expected paths to reject extras and skip the
output-tree block when the output root is absent. A separate check must establish
existence for all 77 required paths. Final conditional-route reconciliation may
record \code{requested} without \code{ran}. Published-table projection comparison
drops TeX files without projected numeric tokens and ignores PDFs, SVGs, CSVs,
RDS files, Markdown files, the descriptive wrapper, and all other non-table
artifacts.

The following claims must remain separate:
\begin{enumerate}
  \item the artifact manifest has valid schema, naming, and path topology,
    every producer and consumer code is used, and every listed producer path
    resolves;
  \item a typed CSV artifact passed its typed CSV round-trip check;
  \item an RDS artifact passed exact object RDS readback;
  \item a published numerical TeX table passed its projection comparison;
  \item every required artifact and every conditional artifact whose producer
    ran exists and is fresh;
  \item every renderable publication output renders correctly.
\end{enumerate}
Dedicated checks cover manifest topology, typed CSV round trips, exact object
RDS readback, and published numerical table projections. Artifact completeness,
freshness, and render acceptance require a separate criterion. Current code
lacks a single implementation of that criterion; final conditional-route
reconciliation covers only conditional routes.

\subsection{Static source trace versus a completed pipeline run}

This trace establishes agreement between the documentation and current source.
Runtime evidence is required for initialized outputs, cache validity, external
data, LAD availability, numerical convergence, and production completion. This
review traced source, mapped the analysis graph, compiled only this standalone
TeX file, and checked each claim against code.

\clearpage
\section{Terminology and code-name reference}
\label{sec:terminology}

This reference preserves the document's one-concept/one-word rule and keeps
literal source names visible. Its four parts distinguish scientific concepts,
operational concepts, technical language, and similar names that denote
different concepts.

\small
\begin{longtable}{@{}>{\raggedright\arraybackslash}p{0.20\textwidth}
  >{\raggedright\arraybackslash}p{0.32\textwidth}
  >{\raggedright\arraybackslash}p{0.41\textwidth}@{}}
\caption{Canonical scientific terms and their current code names.
\label{tab:scientific-terms}}\\
\toprule
Canonical term & Code names & Why the names denote one concept\\
\midrule
\endfirsthead
\toprule
Canonical term & Code names & Why the names denote one concept\\
\midrule
\endhead
quarterly real consumption growth & \code{gr1_pcecc96}; source series
\code{PCECC96} & The data stage transforms \code{PCECC96} into the one quarterly
series that becomes the mean-equation outcome.\\
SDF-news panel & \code{sdf_news}; comments that say SDF news, news SDF, or SDF
innovations & These names refer to the same realized SDF-news panel produced by
\src{scripts-paper/data_preparation/build_sdf_series.R}{sdf_news}.\\
expected-SDF panel & \code{expected_sdf} & The expected-SDF panel is
contemporaneously dated and produced with the SDF-news panel.\\
lagged expected-SDF panel & \code{lag_expected_sdf} & The data stage relabels
the expected-SDF panel one quarter forward for predictive alignment. This
operation preserves the expectation concept.\\
expected-SDF PC score block & \code{expected_sdf_pc} & A fresh PCA of the
expected-SDF panel produces this three-score block.\\
lagged expected-SDF PC score block & \code{lag_expected_sdf_pc} & A fresh PCA
of the lagged expected-SDF panel produces this three-score block, whose signs
are aligned with the expected-SDF PC score block. It is not a date relabeling of
that score block.\\
SDF-news PC score block & \code{sdf_news_pc} & A fresh PCA of the SDF-news
panel produces this three-score block.\\
three SDF-PC score blocks & \code{expected_sdf_pc},
\code{lag_expected_sdf_pc}, and \code{sdf_news_pc} & This collective term means
exactly the expected-SDF PC score block, lagged expected-SDF PC score block, and
SDF-news PC score block. It does not include the nominal asset-return PCs.\\
expected-SDF PC score & columns \code{expected_sdf_pc1}--
\code{expected_sdf_pc3} & Each is one column of the expected-SDF PC score block.\\
lagged expected-SDF PC score & columns \code{lag_expected_sdf_pc1}--
\code{lag_expected_sdf_pc3} & Each is one column of the lagged expected-SDF PC
score block.\\
SDF-news PC score & columns \code{sdf_news_pc1}--\code{sdf_news_pc3}; generic
code comments saying $W_2$ component or news component & The current paper
application supplies the three columns of the SDF-news PC score block as
$Y_2$. Generic identification code calls a $W_2$ column a component; this
document calls each such column in the current application an SDF-news PC score.\\
nominal asset-return PCs & \code{lag_asset_return_pc}, \code{l.pc1}--\code{l.pc4},
and configuration fields containing \code{return_pc} & The data stage reads one
bundled four-PC panel of nominal financial-asset returns and relabels its dates.
These PCs are never yield PCs.\\
five-year-yield volatility instrument & \code{yield_vol[["y60_vol"]]},
\code{y60_vol}, \code{z_col}, and mathematical $Z$ & The yield-volatility panel
contains many maturities; configuration selects its five-year column as the one
identification instrument, which the estimator demeans.\\
published mean specification B & \code{"B"},
\code{set_id_mean_eq_by_spec[["B"]]}, and \code{set_id_mean_eq} &
\src{scripts-paper/mean_equation/estimate_identified_set.R}{set_id_mean_eq}
binds the published entry selected by \code{PAPER_SPEC_PLAN}; that entry is B.\\
diagnostic mean specification A & \code{"A"} and the nonpublished entry of
\code{set_id_mean_eq_by_spec} & The same mean estimator runs with the A
restriction and retains A only for the paired B-versus-A mean-specification
comparison.\\
mean news coefficient $\theta$ & code variables \code{theta}, \code{b},
\code{b_N}, \code{b_point}, and \code{theta_table} in mean-identification and
log-variance set map input contexts & In those contexts the mean news coefficient
is the three-vector multiplying $W_2$ in
$e_t(\theta)=W_{1t}-W_{2t}^{\top}\theta$. The log-variance code commonly calls
the same input \code{b} to distinguish it from its own coefficient vector.\\
residual-scale coefficient $\eta$ & code variables \code{theta},
\code{theta_hat}, \code{coef}, and occasional \code{eta} & These are the five
estimator-specific residual-scale coefficients multiplying $R_t$. PPML and the
Harvey estimator target log conditional variance; the two-step log-OLS
benchmark targets a mean-log projection; LAD targets a conditional-median log
scale. The prose uses $\eta$ to keep these coefficients distinct from the mean
news coefficient $\theta$, even when estimator code calls them \code{theta}.\\
relaxation value $\tau$ & \code{tau}; code comments saying slack, display slack,
or critical slack & All enter the same relaxed heteroskedasticity restrictions
and the same quadratic-system builder. Numerical tolerances are separate.\\
mean identified set $\Theta(\tau)$ & \code{set_id_mean_eq}, entries of
\code{set_tables}, comments saying news set, and per-$\tau$ quadratic systems &
These names denote representations or evaluations of the same feasible set of mean news
coefficients at a fixed relaxation value. Bootstrap inference objects such as
\code{set_id_boot} are not the mean identified set itself.\\
tau-zero point & Lewbel point, point benchmark, \code{point0},
\code{theta_table$point}, and \code{b_point} & The rank-and-consistency solve at
$\tau=0$ produces one mean news coefficient; downstream names pass that same
point into log-variance estimators. An estimator's fit at that mean point is a
separate tau-zero log-variance fit.\\
estimated boundedness frontier $\tau^\star$ & critical slack, \code{tau_star},
and bounded-to-unbounded transition & These names denote the same estimated
transition in the mean identified set as the common relaxation value increases.
When uncapped, the search locates it from certified bounded and unbounded evaluations;
a capped value is the sweep maximum, not a certified transition.\\
endpoint-attaining point & \code{arg_lower}, \code{arg_upper}, \code{arg}, and
comments saying attaining point or endpoint witness in log-variance engine
results & Each is a feasible mean news coefficient at which one reported
log-variance endpoint is attained.
The joint-null witness is a different diagnostic object.\\
stored endpoint status & endpoint state, status word, side status,
\code{PAPER_ENDPOINT_STATUS}, and \code{point_status} & All stored endpoint
classifications use the vocabulary \code{bounded}, \code{unbounded},
\code{unreliable}, and \code{failed}. Ordinary solvers emit \code{bounded},
\code{unbounded}, and \code{unreliable}; collection code can add
\code{failed}.\\
shared log-variance sample & common volatility sample, frozen sample, frozen
inputs, and \code{log_var_eq$sample_contract} &
\src{scripts-paper/log_variance/estimators/log_ols/run.R}{log_var_eq} constructs
one quarter-aligned sample and design reused by all log-variance estimators.\\
log-variance set engine & generic engine, shared set engine, and functions under
\file{scripts-paper/log_variance/engine/} & PPML, Harvey, and the other mapped
estimators call this one implementation for grids, endpoint searches,
certificates, fit-evaluation budgets, caches, and result schemas.\\
primary full-sample set map & primary map, full-sample map, and
\code{primary_results} & For each estimator, these names denote the first
set-engine evaluation over every displayed relaxation value.\\
second search audit & PPML coverage gate/run/audit and Harvey sensitivity
gate/audit & Both call
\src{scripts-paper/log_variance/estimators/audit_orchestration.R}
{logvar_audit_display_taus}. Both reconcile through
\src{scripts-paper/log_variance/estimators/ppml/coverage.R}
{logvar_ppml_apply_coverage}. Their common role is a rerun of every displayed
relaxation value using one fresh set-engine cache, five scan-derived starts per
coefficient and endpoint side, and a fresh 40,000-fit-evaluation budget per
displayed relaxation value.\\
PPML grid-coverage check & PPML coverage grid, Morton selector, and
\code{ppml_coverage$grid_cap} & The PPML grid-coverage check is the extra
component inside PPML's second search audit. It raises the selected-grid cap
from 4,000 to 8,000 and replaces deterministic lattice coarsening with Morton
selection. Harvey's second search audit keeps the primary 4,000-row grid rule.\\
fit-evaluation budget & \code{fit_budget}, \code{max_fit_evals}, and
\code{sensitivity_fit_budget} & Each name limits estimator-fit claims within one
set-engine search. Separate controls set the selected-grid cap and bootstrap
replication count.\\
selected-grid cap & \code{grid_cap}, \code{max_grid_points}, and prose saying
grid cap & Each limits how many feasible grid rows the set engine retains for a
scan. It does not count fit evaluations.\\
set-engine cache & \code{cache}, \code{cache$store}, and comments saying fresh
cache in the log-variance set engine & A set-engine cache is an in-memory memo
of estimator fits within one set map or second search audit. A fresh set-engine cache is a new
empty environment, not a disk cache or the unified bootstrap cache artifact.\\
attained mean coefficient range & code and filenames saying mean bound,
\code{mean_eq_bounds_tau}, \code{set_lower}, and \code{set_upper} & When both
endpoint statuses reduce to \code{bounded}, the two finite endpoint values form
a certified attained coordinate or affine-functional range. Otherwise the row
is open or uncertified. The term \emph{attained} preserves the source caveat
that multistart search does not prove globally sharp extrema.\\
bounds-by-relaxation-value figure & code and artifact names containing
\code{bounds_tau}, \code{render_bounds_by_tau}, or bounds-by-tau & These are
coefficient-endpoint figures indexed by the identification relaxation value
$\tau$. They are unrelated to the SDF approximation-error variance bounds.\\
exogenous-news OLS benchmark & \code{ols_mean_eq}; comments saying same-sample
OLS or naive reference & This regression treats the three SDF-news PC scores as
exogenous.
It is a comparison benchmark, not the structural mean identified-set estimate.\\
structural design coefficient $\beta_1(\theta)$ & \code{beta1},
\code{beta1_point}, \code{beta1_table}, and
\code{recover_structural_coefficients} & The exact aligned-sample identity is
$\beta_1(\theta)=\beta_1^R-(\beta_2^R)^\top\theta$. The reduced-form coefficient
$\beta_1^R$ is a separate input.\\
attained variance-share range & set-valued share, \code{var_share}, and set
columns in the variance-share table & These are attained numerical ranges of the
same variance-share maps over the mean identified set. They are not certified
global extrema.\\
two-step log-OLS benchmark & log-OLS foundation, mean-log benchmark,
\code{log_var_eq}, and registry key \code{logols} & One analysis stage both
freezes the shared log-variance sample and maps
$\log(e_t(\theta)^2)$ by OLS. It is a benchmark outside the primary
log-variance estimator registry.\\
PPML quasi-Poisson log-link QMLE & PPML, \code{log_var_eq_ppml}, registry key
\code{ppml}, and target names containing \code{theta_var} & These names denote
one conditional-variance coefficient estimator and its set map.\\
Harvey Gaussian multiplicative-variance QMLE & Harvey estimator,
\code{log_var_eq_harvey}, registry key \code{harvey}, and target names containing
\code{theta_var_gaussian} & These names denote one Gaussian
conditional-variance coefficient estimator and its set map. The Harvey
heteroskedasticity test is unrelated.\\
dependency-gated LAD conditional-median map & LAD, \code{log_var_eq_lad}, registry key
\code{lad}, and target \code{theta_median} & These names denote the
dependency-gated median-regression map for $2\log|e_t(\theta)|$. Its fitted
figure is a conditional-median absolute-residual envelope, not the conditional
residual volatility targeted by PPML and the Harvey estimator.\\
joint-null witness diagnostic & joint-null stage, witness diagnostic,
\code{log_var_eq_joint_null}, and artifact IDs beginning \code{joint_null} & One
search asks whether a feasible mean news coefficient is compatible with jointly
zero nominal asset-return-PC slopes. It is not a hypothesis test.\\
joint-GMM replication diagnostic & joint-GMM branch, graph-replication
diagnostic, \code{log_var_eq_joint_gmm}, and artifact IDs beginning
\code{joint_gmm} & The current branch checks graph, algebra, and software
invariance. It is not an overidentification test or an added restriction.\\
residual-dynamics gate & dynamics gate, Ljung--Box gate, and artifact ID
\code{dynamics_gate} & One predeclared lag-4 screen produces the
residual-dynamics gate record read by the EGARCH-X decision router. It does not
select between PPML and the Harvey estimator.\\
heteroskedasticity diagnostic battery & heteroskedasticity evidence, battery,
tests, and the three-panel heteroskedasticity table & These phrases refer to the
late recomputed collection of named diagnostics. The individual tests remain
distinct, and the Harvey heteroskedasticity test is not the Harvey estimator.\\
EGARCH-X decision router & EGARCH routing, route result, dynamic workstream, and
functions under \file{scripts-paper/log_variance/extensions/egarch/} & The
current code validates the residual-dynamics gate and routes the committed
decision. It contains no EGARCH-X artifact producer.\\
SDF approximation-error variance bounds & variance-bound stage,
\code{variance_bounds_df}, and the news/expected-SDF bound artifact IDs & These
are maturity-specific plug-in estimates of population approximation-error
bounds. They are unrelated to coefficient endpoints over the mean identified
set.\\
first-order-cancelled q-arm & q-arm and report strings saying no-Taylor q-arm &
The q construction retains the full first-order-cancelled remainder. Envelope
and component arms are different formulas whose retained leading error is
fourth order.\\
tau-zero bootstrap point/MAD statistic & bootstrap t statistic,
\code{point_t_statistic}, and \code{point_t} & One statistic divides the
full-sample point by the MAD of bounded point draws when the tau-zero reporting
gate passes. Otherwise the statistic and p-values are unavailable. The
associated bootstrap absolute-deviation p-value is not studentized, even though
table-facing code uses the t-statistic label.\\
\TauP{} confidence interval & \code{target_p_critical}, \code{ci_lower}, and
\code{ci_upper} & When the endpoint regularity gate passes and the full-sample
status pattern supports a finite side, the Target P confidence interval is the published pointwise-coverage
interval, uniform over the relevant scalar identified interval. Otherwise the
interval is suppressed or the requested-relaxation-value gate stops
publication. Mean-coefficient intervals apply to scalar projections of the mean
identified set. The confidence limits are not identified-set endpoints.\\
\TauS{} containment diagnostic & \code{target_s_critical} and \code{c_s} & This
diagnostic calibrates containment of the full identified set. It is not the
published \TauP{} target.\\
endpoint regularity gate & minimum valid-draw share, status-stability share,
positive MAD scale, and common-pair count & These code fields implement one
shared publication gate for mean and log-variance endpoint inference. It is
unrelated to routing or dependency gates.\\
\bottomrule
\end{longtable}
\normalsize

Table~\ref{tab:scientific-terms} keeps code matching explicit without allowing
source aliases to proliferate in the prose.

\small
\begin{longtable}{@{}>{\raggedright\arraybackslash}p{0.22\textwidth}
  >{\raggedright\arraybackslash}p{0.31\textwidth}
  >{\raggedright\arraybackslash}p{0.40\textwidth}@{}}
\caption{Canonical operational terms and their current code names.
\label{tab:operational-terms}}\\
\toprule
Canonical term & Code names & Meaning and mapping rule\\
\midrule
\endfirsthead
\toprule
Canonical term & Code names & Meaning and mapping rule\\
\midrule
\endhead
source module & a file passed to \code{paper_source_once()} & A source module is
an R file evaluated into the shared environment. It may define functions
without performing an analysis stage.\\
analysis stage & stage, branch, or top-level driver in comments & An analysis
stage is one substantive unit of data construction, estimation, inference,
diagnosis, or publication. One stage may source several modules.\\
always attempted & unconditional top-level source or call & The runner reaches
the statement whenever every preceding statement succeeds. This status does not
claim that the statement completes.\\
selected path & current checked-in source plus the default launch environment &
This is the branch chosen by a plain invocation. External prerequisites or a
runtime-selected outcome may still stop it.\\
alternative path & nondefault environment input or tracked source decision & The
branch is implemented but is not chosen by a plain invocation. Its exact trigger
is stated where the path is documented.\\
source-selected & checked-in contract, plan, or decision record & The source
chooses between the selected path and an alternative path.\\
launch-selected & \code{Sys.getenv()} branch & An environment input chooses
between the selected path and an alternative path. The document always states
the unset-input path.\\
runtime-selected outcome & cache validation, external-data resolution,
diagnostic verdict, or numerical status & The run chooses the outcome from
current state or data. Static source inspection states the rule, not the
realized outcome.\\
available but not called & inactive, offline, test-support, or
source-loaded but uncalled code & The functionality exists, but the current
runner has no production call that executes it.\\
not implemented & reserved schema field, route result, or artifact record without
a substantive solver or producer & The interface reserves the concept, but the
implementation does not exist.\\
separate entrypoint & reset, test, validation, quality, or direct package command
or API call & Another entrypoint owns the functionality. The paper runner does
not invoke it.\\
artifact producer & producer, writer, or renderer when it writes a manifested
path & An artifact producer is the file, set of files, or directory named by an
artifact record's \code{producer} field. This document uses the term
\emph{renderer} only for a rendering function. The manifest role remains
artifact producer.\\
source-once registry & \code{.paper_source_registry}; source-state registry &
The source-once protocol distinguishes an absent binding, \code{FALSE} while
loading, and \code{TRUE} after loading for each normalized source path.\\
log-variance estimator registry & \code{PAPER_LOGVAR_PRIMARY_ESTIMATORS},
\code{PAPER_LOGVAR_ESTIMATOR_EXTENSIONS}, and \code{PAPER_LOGVAR_ESTIMATORS} &
The combined registry stores estimator identities, dependencies, capabilities,
display metadata, artifact IDs, and budget policies. Membership in the two
source lists determines primary versus extension status.\\
bounds-rendering registry & \code{logvar_bounds_tau_registry}; comments saying
figure registry or bounds registry & This runtime list stores estimator objects,
set map results, search seeds, and engine options for the bounds-by-relaxation-value
artifact producer.\\
artifact manifest; artifact record; artifact ID & \code{artifact_manifest};
comments saying manifest row or file record & The data frame is the artifact
manifest. Each row is an artifact record. The row's \code{id} field is the
artifact ID. Family and variant are separate lookup fields.\\
artifact group & \code{artifact_manifest$group} and
\code{PAPER_ARTIFACT_GROUPS} & The group selects the output subdirectory. It is
not an artifact lifecycle status or an artifact lookup family.\\
artifact lifecycle status & \code{artifact_manifest$status},
\code{required}, \code{conditional_lad}, and \code{conditional_egarch} & The
status classifies a record as required or as part of a conditional lifecycle
route. Conditional statuses drive the startup cleanup of conditional artifacts
and final reconciliation. It
is unrelated to stored endpoint status.\\
artifact lookup family/variant & \code{artifact_manifest$family},
\code{artifact_manifest$variant}, and comments saying family lookup & This
nullable coordinate selects one artifact record. It is unrelated to a
resample-index family.\\
LaTeX publication registry; publication triple &
\code{artifact_latex_publications} and LaTeX publication registry & Each runtime
row is one publication triple: fragment, standalone TeX, and standalone PDF
artifact IDs. \code{.latex_publication_specs} is temporary construction input,
not another name for the runtime registry.\\
unified bootstrap stage & single shared stage, unified stage, and
\code{run_bootstrap_stage} & One stage constructs both resample-index families,
then computes or reuses the unified result payload. When the stage computes the
payload, it evaluates published mean and log-variance result draws. It owns the
unified bootstrap cache artifact.\\
primary resample-index family & \code{bootstrap_primary_family}, primary
family, and resample-index family & The primary resample-index family is the immutable record containing the
ordered circular-moving-block index draws and their design, seed, RNG-state, and
hash metadata. \emph{Primary indices} means its \code{indices} field. Published
mean specification B and diagnostic mean specification A use that same field.
Estimator-result draws are not a resample-index family.\\
volatility-sensitivity resample-index family & sensitivity family,
doubled-block family, and volatility sensitivity family & The volatility-sensitivity resample-index family is the immutable
record containing the ordered index draws with block length twice that of the
primary resample-index family and the same kinds of protocol metadata. Its
\code{indices} field drives volatility-only re-estimation; it does not rescale a
standard error.\\
unified bootstrap cache artifact & unified state artifact, bootstrap cache, and
artifact ID \code{bootstrap_stage_draws} & The RDS stores the unified stage's
result collections, failure counts, volatility anchor, and provenance. Both
resample-index families are rebuilt before cache validation. Provenance stores
their design and identity hashes, not their family objects or index vectors.\\
bootstrap cache payload; bootstrap cache schema 4 &
\code{BOOTSTRAP_STAGE_CACHE_FIELDS}, \code{BOOTSTRAP_STAGE_CACHE_SCHEMA}, and
\code{bootstrap_stage_cache_payload} & The payload is the seven-field R object
stored in the unified bootstrap cache artifact. Schema 4 is its validation
contract, not a route or artifact manifest schema.\\
bootstrap provenance record; semantic cache freshness & provenance, semantic
freshness, cache identity, and cache reuse validity & The provenance record
stores the reuse-gating design, input, code, runtime, and resample-index hashes.
Semantic freshness means equality under those validators, not modification-time
freshness or cryptographic authentication.\\
bootstrap draw-code manifest; bootstrap presentation manifest &
\code{bootstrap_stage_code_manifest()}, \code{code_sha},
\code{bootstrap_stage_presentation_manifest()}, and \code{presentation_sha} &
Both enumerate source paths for hashing. Only the bootstrap draw-code manifest
gates cache reuse; the presentation digest is audit-only.\\
conditional-route state record & final route record, route-status object, and
the result of \code{build_conditional_route_status()} & This in-memory record
reconciles route decisions, requests, producer execution, and conditional-artifact
presence. It carries audits of the startup cleanup of conditional artifacts
unchanged for provenance and is
persisted as the \code{conditional_route_status} artifact.\\
route decision; requested; ran; artifact present & fields in the
conditional-route state record & These four fields describe different points in
the routing chain: decision, permission, producer execution, and path existence.
A requested route may validly have \code{ran=FALSE}.\\
startup cleanup of conditional artifacts; full pipeline-state reset; final
LaTeX-sidecar cleanup & \code{cleanup_conditional_artifacts},
\code{reset_pipeline_state}, and \code{clean_latex_sidecars} & These operations
share lifecycle infrastructure but have different scopes and call sites. The
runner performs the first and last, not the full reset.\\
topology checks & topology audit, topology tests, and topology check &
These names refer to the suite that parses and style-checks the R tree, checks
source references and source-once behavior, validates artifact manifest schema
and path formulas, verifies that listed producer paths resolve, and rejects
sidecars, root files, and unmanifested extras when an output tree exists. They
do not establish artifact completeness or producer ownership.\\
typed CSV round-trip check; exact object RDS readback & typed writer readback,
CSV readback, or exact-RDS readback & A typed CSV writer performs a
tolerance-aware data-frame round trip. An exact-RDS writer requires
\code{identical()} object equality after readback. Neither establishes byte
identity.\\
published-table projection comparison & table comparator, table comparison, and
numerical TeX comparison & The validation code compares projected table paths,
numerical coordinates, token counts, numerical values under a precision-aware
rounding rule, and attached stars across table trees. It does not require equal
textual precision.\\
criterion for artifact completeness and render acceptance & completeness pass,
publication acceptance, or render acceptance in descriptive comments & This
additional criterion would require every required artifact and every produced
conditional artifact whose producer ran to exist and be fresh, and every
renderable output to render correctly. Current code has not implemented it as
one check.\\
static source trace & source inspection, code trace, and static trace & A static
source trace follows current source without executing
the pipeline or treating generated outputs as evidence.\\
\bottomrule
\end{longtable}
\normalsize

Table~\ref{tab:operational-terms} shows why shared code words such as
\emph{state}, \emph{family}, and \emph{manifest} do not imply one operational
object.

\small
\begin{longtable}{@{}>{\raggedright\arraybackslash}p{0.20\textwidth}
  >{\raggedright\arraybackslash}p{0.37\textwidth}
  >{\raggedright\arraybackslash}p{0.36\textwidth}@{}}
\caption{Plain-English decoder for technical abbreviations and numerical-search
terms. The source column identifies the current pipeline use, not a general
textbook definition.\label{tab:technical-terms}}\\
\toprule
Term & Meaning in this document & Current source authority\\
\midrule
\endfirsthead
\toprule
Term & Meaning in this document & Current source authority\\
\midrule
\endhead
API & Application programming interface: a callable interface between code
layers or modules. APIs in this document include the package identification API
and the paper-owned artifact query API. & \file{R/};
\file{scripts-paper/config/artifacts.R};
\file{scripts-paper/config/artifact_registry_queries.R}\\
SDF & Stochastic discount factor. The SDF-news panel is its realized news
component; the expected-SDF panel is its conditional expectation. &
\file{scripts-paper/data_preparation/build_sdf_series.R}\\
ACM & Adrian--Crump--Moench yield and term-premium inputs. The validated
quarterly ACM loader supplies the maturity panels used by the data and
approximation-error-bound stages. &
\file{scripts-paper/support/data/acm_inputs.R}\\
FRED & Federal Reserve Economic Data, the external source used for the real
consumption series. &
\file{scripts-paper/data_preparation/fred_download_patch.R};
\file{scripts-paper/data_preparation/build_consumption_growth.R}\\
PC; PCA & A PC is a principal-component score. PCA is the principal-component
analysis that computes scores. The pipeline loads nominal asset-return PC scores
and computes the three SDF-PC score blocks. &
\file{scripts-paper/data_preparation/build_asset_return_pcs.R};
\file{scripts-paper/data_preparation/build_sdf_pcs.R}\\
OLS & Ordinary least squares. The exogenous-news OLS benchmark and the two-step
log-OLS benchmark are different regressions with different outcomes and roles. &
\file{scripts-paper/mean_equation/fit_ols.R};
\file{scripts-paper/log_variance/estimators/log_ols/estimator.R}\\
QMLE & Quasi-maximum-likelihood estimator. In this pipeline, PPML and the Harvey
estimator are two different QMLEs for the conditional-variance map. &
\file{scripts-paper/log_variance/estimators/ppml/fit.R};
\file{scripts-paper/log_variance/estimators/harvey/likelihood.R}\\
PPML & Poisson pseudo-maximum likelihood. The current implementation is the
quasi-Poisson log-link QMLE fitted with \code{glm.fit()}. It is one estimator,
not a generic name for every log-variance set map. &
\file{scripts-paper/log_variance/estimators/ppml/fit.R}\\
LAD & Least absolute deviations, implemented here as median quantile regression
for the conditional-median log scale. The \code{tau=0.5} argument is a quantile
probability, not the identification relaxation value. &
\file{scripts-paper/log_variance/estimators/lad/estimator.R}\\
GMM & Generalized method of moments. The current joint-GMM replication
diagnostic checks graph, algebra, and software invariance; it is not an added
estimator or identifying restriction. &
\file{scripts-paper/log_variance/diagnostics/joint_gmm/pipeline_driver.R}\\
EGARCH-X & Exponential generalized autoregressive conditional
heteroskedasticity with exogenous predictors. The current runner contains its
decision router but no EGARCH-X artifact producer. &
\file{scripts-paper/log_variance/extensions/egarch/decision_core.R};
\file{scripts-paper/log_variance/extensions/egarch/run_route.R}\\
circular MBB & Circular moving-block bootstrap. A resample-index family stores
ordered index draws made from circular blocks; it is not a collection of
coefficient estimates. &
\file{scripts-paper/support/statistics/bootstrap_and_stationarity.R};
\file{scripts-paper/support/statistics/mbb_index_family.R}\\
RNG & Random-number generator. The resampling protocol fixes the generator,
normal-generator, and sampling kinds as Mersenne-Twister, Inversion, and
Rejection, and records their identity in provenance. &
\file{scripts-paper/support/statistics/mbb_protocol_authority.R};
\file{scripts-paper/support/inference/bootstrap_stage_provenance.R}\\
HAC & Heteroskedasticity-and-autocorrelation-consistent covariance estimator.
The reporting rule selects a four-lag Bartlett version for PPML and the Harvey
estimator. &
\file{scripts-paper/log_variance/inference/standard_error_estimators.R};
\file{scripts-paper/config/reporting.R}\\
SE; HC0; HC1 & SE means standard error. HC0 and HC1 are two
heteroskedasticity-consistent PPML covariance variants; the current reporting
policy selects HAC instead. &
\file{scripts-paper/log_variance/estimators/ppml/standard_errors.R};
\file{scripts-paper/config/reporting.R}\\
OPG & Outer product of gradients, one Harvey estimator covariance variant. It
is not the reporting covariance selected by the current policy. &
\file{scripts-paper/log_variance/estimators/harvey/standard_errors.R}\\
MAD & Median absolute deviation. The code stores this normal-consistent robust
bootstrap scale as \code{se}/\code{point_se} and uses it in the tau-zero
bootstrap point/MAD statistic. It is not the sample standard deviation, and the
common scale cancels from the associated bootstrap p-value. &
\file{scripts-paper/support/inference_post/endpoint_point_statistic.R}\\
SLSQP & Sequential least-squares quadratic programming, the local constrained
optimizer used in mean-geometry refinement and several endpoint searches. A
successful local solve does not by itself prove a global endpoint. &
\file{scripts-paper/support/identification/profile_solver_core.R};
\file{scripts-paper/log_variance/core/endpoint_polish.R}\\
COBYLA & Constrained Optimization BY Linear Approximations, the derivative-free
local constrained optimizer used for nonsmooth objectives or maps without a
gradient. LAD endpoint polishing follows this path. Like SLSQP, it does not by
itself prove a global endpoint. &
\file{scripts-paper/log_variance/core/endpoint_polish.R};
\file{scripts-paper/log_variance/estimators/lad/estimator.R}\\
Morton selection & A deterministic, space-filling ordering used only by the
PPML grid-coverage check to select at most 8,000 feasible grid rows. &
\file{scripts-paper/log_variance/estimators/ppml/pilot_and_grid.R}\\
Lipschitz branch-and-bound & The certified one-dimensional global search used
by Target P to bound the worst truth location over a scalar identified interval.
It stops when the gap between its lower and upper reported-value bounds is
within tolerance. It is distinct from local SLSQP or COBYLA endpoint polishing.
& \file{scripts-paper/support/inference_post/endpoint_targets.R}\\
warm start; scan-derived start pool; endpoint polishing; cold-start endpoint
comparison & A warm start passes the previous successful fit's warm-start state
to the next grid fit along the traversal. The scan-derived start pools retain
separated feasible candidates for each coefficient and endpoint
side. Endpoint polishing launches local constrained optimization from chosen
candidates. The cold-start endpoint comparison independently refits an accepted
endpoint without the set-engine cache. These are four distinct search steps. &
\file{scripts-paper/log_variance/engine/grid_scan.R};
\file{scripts-paper/log_variance/engine/endpoints.R};
\file{scripts-paper/log_variance/engine/results.R}\\
numerical certificate & The umbrella term for context-specific,
solver-independent acceptance evidence required before the code labels an
endpoint \code{bounded} or a direction \code{unbounded}. Mean bounds,
nonlinear residual-scale endpoints, log-OLS open directions, and Harvey
nonexistence use different certificates. Solver convergence or a finite number
alone is not a certificate. &
\file{scripts-paper/support/identification/profile_solver_core.R};
\file{scripts-paper/log_variance/core/endpoint_polish.R};
\file{scripts-paper/log_variance/core/residual_map.R};
\file{scripts-paper/log_variance/estimators/harvey/recession_certificate.R}\\
Hamilton rounding & Largest-remainder rounding used only for fixed-coefficient
variance-share columns so displayed components add to displayed block totals. &
\file{scripts-paper/mean_equation/variance_shares/render_variance_share_table.R}\\
LM; ARCH(1) & LM means Lagrange-multiplier statistic. ARCH(1) is a first-order
autoregressive-conditional-heteroskedasticity diagnostic. They are distinct
members of the heteroskedasticity diagnostic battery. &
\file{scripts-paper/mean_equation/diagnostics/heteroskedasticity/}\\
BLAS; LAPACK & Numerical linear-algebra libraries. Their runtime identities enter
the bootstrap provenance hash because they may affect numerical computation. &
\file{scripts-paper/support/statistics/boot_freshness.R};
\file{scripts-paper/support/inference/bootstrap_stage_provenance.R}\\
SD units & Standard-deviation units, used as one figure scaling alongside raw
coefficient units. &
\file{scripts-paper/mean_equation/figures/prepare_region_geometry.R}\\
CSV; RDS; TeX; PDF; SVG & Distinct artifact formats: comma-separated values,
one serialized R object, TeX source, Portable Document Format, and Scalable
Vector Graphics. Artifact IDs make them addressable; the selected writer
determines typed CSV or exact-object RDS readback behavior. &
\file{scripts-paper/support/artifacts/typed_artifacts.R};
\file{scripts-paper/config/artifact_latex.R}\\
\bottomrule
\end{longtable}
\normalsize

Table~\ref{tab:technical-terms} separates statistical abbreviations from search
controls: an optimizer name or convergence message never substitutes for a
numerical certificate.

\small
\begin{longtable}{@{}>{\raggedright\arraybackslash}p{0.31\textwidth}
  >{\raggedright\arraybackslash}p{0.62\textwidth}@{}}
\caption{Similar names that denote different concepts.\label{tab:false-friends}}\\
\toprule
Keep separate & Reason\\
\midrule
\endfirsthead
\toprule
Keep separate & Reason\\
\midrule
\endhead
Source-once registry, log-variance estimator registry, bounds-rendering
registry, artifact manifest, and LaTeX publication registry & The source-once
registry stores loading state. The log-variance estimator registry stores
estimator specifications. The bounds-rendering registry stores runtime rendering
inputs. The artifact manifest stores artifact metadata records. The LaTeX
publication registry stores publication triples.\\
Relaxation value and numerical tolerance & A relaxation value changes the
scientific identified set. A numerical tolerance controls a comparison at a
fixed scientific problem.\\
Identification relaxation value $\tau$ and LAD quantile probability
\code{tau=0.5} & The identification relaxation value indexes the mean identified
set. The LAD quantile probability is the
median-regression probability passed to \code{quantreg} and never changes the
mean identified set.\\
$W_1$, $e_t(\theta)$, and $u_t$ & $W_1$ is quarterly real consumption growth
residualized on $X$. The residual $e_t(\theta)=W_{1t}-W_{2t}^{\top}\theta$
varies over the mean identified set. The residual $u_t$ comes from the
exogenous-news OLS benchmark.\\
Endpoint-attaining point and joint-null witness & An endpoint-attaining point
attains one residual-scale coefficient endpoint. A joint-null witness is a
feasible mean news coefficient compatible with zero nominal asset-return-PC
slopes in a diagnostic.\\
Two-step log-OLS benchmark, PPML, Harvey estimator, and LAD & The two-step
log-OLS benchmark targets a mean-log projection. PPML and the Harvey estimator
target conditional variance. LAD targets conditional median scale. The four
estimators have different criteria.\\
Harvey estimator and Harvey heteroskedasticity test & They are unrelated methods
that share an eponym. Every diagnostic occurrence keeps the word
\emph{heteroskedasticity test}.\\
Second search audit and PPML grid-coverage check & The second search audit is the
common rerun with a fresh set-engine cache used by PPML and the Harvey estimator.
The PPML grid-coverage check is the PPML-only change to an 8,000-row
Morton-selected grid inside that audit.\\
Second search audits for PPML and the Harvey estimator, and LAD verify-or-rerun
check & The PPML and Harvey second search audits reconcile endpoints and may
demote a side. The LAD verify-or-rerun check verifies or reruns a row and does
not patch endpoints.\\
PPML and Harvey estimator fitted conditional-residual volatility, and LAD fitted
conditional-median scale & Fitted conditional-residual volatility is conditional
residual standard deviation. Fitted conditional-median scale is conditional
median absolute-residual magnitude. The shared artifact lookup family does not
make the estimands equal.\\
Residual-dynamics gate, EGARCH-X decision router, and EGARCH-X artifact producer &
The residual-dynamics gate supplies a diagnostic verdict. The EGARCH-X decision
router supplies a routing decision. An EGARCH-X artifact producer would estimate
the model; the current runner has no such producer.\\
Coefficient bounds by relaxation value and SDF approximation-error variance
bounds & Coefficient bounds by relaxation value are endpoints over the mean
identified set. SDF approximation-error variance bounds are maturity-specific
plug-in estimates of population bounds.\\
Identified-set endpoints and \TauP{} confidence limits & Identified-set endpoints
describe the parameter set. Target P confidence limits quantify sampling
uncertainty about those endpoints.\\
\TauP{} and \TauS{} & \TauP{} targets pointwise coverage uniformly over the
relevant scalar identified interval and is published. \TauS{} targets set
containment and remains a diagnostic.\\
Bootstrap point/MAD statistic and bootstrap absolute-deviation p-value & The
bootstrap point/MAD statistic uses a MAD denominator. The bootstrap
absolute-deviation p-value compares unstudentized absolute deviations because
that common scale cancels.\\
Artifact lookup family, resample-index family, and a resample-index family's
\code{indices} field & An artifact lookup family is a manifest lookup coordinate.
A resample-index family is an immutable resampling record. Its \code{indices}
field contains the ordered index draws.\\
Resample-index family, estimator-result draws, and unified bootstrap cache
artifact & A resample-index family contains resample indices and protocol
metadata. Estimator-result draws contain estimates produced from those indices.
The unified bootstrap cache artifact persists result collections and provenance.
The cache stores family hashes, not family
objects or index vectors.\\
Artifact group, artifact lifecycle status, and artifact lookup family/variant &
An artifact group selects an output directory. An artifact lifecycle status
marks a record as required or conditional. An artifact lookup family/variant
selects a nullable lookup coordinate. The three classifications are
independent.\\
Source-loading state, artifact \code{state} group, stored endpoint status,
conditional-route state record, and unified bootstrap cache artifact &
Source-loading state records module-load progress. The artifact \code{state}
group is an output directory. Stored endpoint status classifies an endpoint.
The conditional-route state record persists route reconciliation. The unified
bootstrap cache artifact persists bootstrap results and provenance. The shared
word \emph{state} in code does not make them one concept.\\
Artifact lifecycle status, execution reachability status, route decision,
requested, producer ran, and artifact present & Artifact lifecycle status
classifies a manifest record. Execution reachability status says whether and how
the runner can reach code. A route decision records the committed branch.
\code{requested} records
permission to execute. \code{ran} records producer execution. Artifact presence
records path existence. These routing facts may differ within one route.\\
Semantic cache freshness, residual-dynamics-record freshness, and artifact
freshness & Semantic cache freshness means provenance equality.
Residual-dynamics-record freshness means matching decision input. Artifact
freshness means acceptance-time file currency.\\
Set-engine cache and unified bootstrap cache artifact & A set-engine cache is an
in-memory estimator-fit memo used within log-variance searches. A unified
bootstrap cache artifact is a
persisted RDS artifact holding unified bootstrap results and provenance.
A fresh set-engine cache and unified bootstrap cache reuse are independent.\\
Artifact manifest, bootstrap draw-code manifest, bootstrap presentation
manifest, and typed-column manifest & The artifact manifest enumerates artifact
metadata records. The bootstrap draw-code manifest enumerates reuse-gating source
paths. The bootstrap presentation manifest enumerates audit-only presentation
paths. The typed-column manifest enumerates CSV column types.\\
Exact object RDS readback and byte identity & The writer requires
\code{identical()} after serialization. Byte identity across R versions or
machines lies outside the contract.\\
Cache reuse validation, conditional-route reconciliation, topology checks,
typed CSV round-trip or exact object RDS readback, published-table projection
comparison, and criterion for artifact completeness and render acceptance &
Cache reuse validation establishes semantic cache freshness. Conditional-route
reconciliation establishes route consistency and conditional-artifact presence.
Topology checks establish source and path topology. Typed CSV round-trip and
exact object RDS readback establish writer-specific readback properties. The
published-table projection comparison establishes equality of selected numeric
table content and attached stars. The criterion for artifact completeness and
render acceptance would establish required-file completeness, freshness, and
successful rendering; the code does not currently implement it as one check.\\
\bottomrule
\end{longtable}
\normalsize

Table~\ref{tab:false-friends} states the distinctions that the canonical
terminology must preserve.

\begin{landscape}
\section{Code-to-document index}
\label{sec:index}

The code-to-document index maps each topic directly to its principal authority.
Transitive helpers remain under the listed directory unless the
index names a more specific source.

\begingroup
\footnotesize
\begin{longtable}{@{}>{\raggedright\arraybackslash}p{0.20\textwidth}
  >{\raggedright\arraybackslash}p{0.66\textwidth}
  >{\raggedright\arraybackslash}p{0.10\textwidth}@{}}
\caption{Code-to-document index. Each row maps one topic to its principal source
authority and to the section that explains the source contract.
\label{tab:code-index}}\\
\toprule
Topic & Principal source authority & Document section\\
\midrule
\endfirsthead
\toprule
Topic & Principal source authority & Document section\\
\midrule
\endhead
Top-level order & \file{scripts-paper/run_pipeline.R} & \ref{sec:runner-ledger}\\
Default route, branch predicates, and reachability status &
\file{scripts-paper/run_pipeline.R}; tracked configuration and decision records &
\ref{sec:status-words}, \ref{sec:current-route}\\
Paths and source-once registry & \file{scripts-paper/config/paths.R} & \ref{sec:runner}\\
Production inputs and prerequisites & \file{scripts-paper/config/paths.R};
\file{scripts-paper/config/analysis_contract.R};
\file{scripts-paper/data_preparation/}; package ACM data APIs under \file{R/};
bootstrap provenance sources & \ref{sec:input-inventory}\\
Supported user controls & \file{scripts-paper/config/analysis.R};
\file{scripts-paper/run_pipeline.R};
\file{scripts-paper/reset_pipeline_state.R};
\file{scripts-paper/config/acquisition.R};
\file{scripts-paper/support/identification/quadratic_system.R};
\file{scripts-paper/config/decisions/lad.dcf} & \ref{sec:user-controls}\\
Source-fixed global control families & \file{scripts-paper/config/};
\file{scripts-paper/support/};
\file{scripts-paper/log_variance/};
\file{scripts-paper/inference/} & \ref{sec:user-controls}\\
Available-but-not-called and separate functionality &
\file{scripts-paper/README.md}; \file{scripts-paper/tests/run_tests.R};
\file{scripts-paper/validation/compare_output_tables.R};
\file{scripts-paper/reset_pipeline_state.R}; \file{docs/quality-check.R} &
\ref{sec:outside-runner}\\
Analysis contract and policies & \file{scripts-paper/config/analysis_contract.R}; \file{scripts-paper/config/analysis.R} & \ref{sec:preflight}\\
Artifact manifest, variants, generated records, and publication mapping & \file{scripts-paper/config/artifacts.R}; \file{scripts-paper/config/artifact_manifest_data.R}; \file{scripts-paper/config/artifact_manifest_variants.R}; \file{scripts-paper/config/artifact_manifest_sweep.R}; \file{scripts-paper/config/artifact_manifest_region.R}; \file{scripts-paper/config/artifact_registry_queries.R}; \file{scripts-paper/config/artifact_latex.R} & \ref{sec:artifacts}\\
Artifact lifecycle and conditional-route state & \file{scripts-paper/config/artifact_lifecycle.R}; \file{scripts-paper/config/artifact_reset.R} & \ref{sec:preflight}, \ref{sec:publication}\\
Data construction & \file{scripts-paper/data_preparation/} & \ref{sec:data}\\
Published mean specification B and diagnostic mean specification A & \file{scripts-paper/mean_equation/estimate_identified_set.R}; \file{scripts-paper/support/identification/identified_set_bootstrap.R} & \ref{sec:specs}\\
Seven centered identification moments and quadratic identified-set system & \file{R/hetid_moments.R}; \file{R/build_general_quadratic_system.R}; \file{scripts-paper/support/identification/quadratic_system.R} & \ref{sec:moments}, \ref{sec:quadratic}\\
Mean identified-set numerical geometry & \file{scripts-paper/support/identification/profile_solver_core.R}; \file{scripts-paper/support/identification/functional_bounds.R}; \file{scripts-paper/support/identification/tau_star.R}; \file{scripts-paper/mean_equation/inference/theta_box_multistart.R} & \ref{sec:quadratic}\\
Attained variance-share ranges & \file{scripts-paper/mean_equation/variance_shares/} & \ref{sec:mean}\\
Log-variance estimator registry and set engine & \file{scripts-paper/config/logvar_estimators.R}; \file{scripts-paper/log_variance/engine/} & \ref{sec:logvar}\\
Two-step log-OLS benchmark, PPML, Harvey estimator, and LAD & \file{scripts-paper/log_variance/estimators/} & \ref{sec:logvar}\\
Diagnostics, residual-dynamics gate, and EGARCH-X decision router & \file{scripts-paper/log_variance/diagnostics/}; \file{scripts-paper/log_variance/extensions/egarch/} & \ref{sec:diagnostics}\\
Unified bootstrap stage & \file{scripts-paper/inference/run_bootstrap_stage.R}; \file{scripts-paper/inference/bootstrap_stage_specs.R}; \file{scripts-paper/inference/bootstrap_stage_draw.R}; \file{scripts-paper/inference/bootstrap_stage_results.R} & \ref{sec:bootstrap}\\
Inference post-processing & \file{scripts-paper/support/inference_post/endpoint_targets.R}; \file{scripts-paper/support/inference_post/endpoint_target_cells.R}; \file{scripts-paper/support/inference_post/endpoint_point_statistic.R} & \ref{sec:bootstrap}\\
Unified bootstrap cache artifact and provenance & \file{scripts-paper/inference/bootstrap_stage_cache.R}; \file{scripts-paper/support/inference/bootstrap_stage_provenance.R}; \file{scripts-paper/support/inference/bootstrap_stage_code_manifest.R}; \file{scripts-paper/support/statistics/boot_cache.R}; \file{scripts-paper/support/statistics/boot_freshness.R} & \ref{sec:cache}\\
Publication artifact producers & \file{scripts-paper/mean_equation/figures/}; \file{scripts-paper/log_variance/figures/}; \file{scripts-paper/log_variance/tables/}; \file{scripts-paper/variance_bounds/}; \file{scripts-paper/reports/} & \ref{sec:publication}\\
Topology checks and published-table projection comparison & \file{scripts-paper/tests/support/check_topology.R}; \file{scripts-paper/tests/support/topology_artifact_checks.R}; \file{scripts-paper/validation/table_comparison.R} & \ref{sec:validation}\\
\bottomrule
\end{longtable}
\endgroup
\end{landscape}

\end{document}
